\begin{document}
\section*{FINAL NEET(UG)-2021 (EXAMINATION)}
\section*{NEET}
(HELD ON SUNDAY $\mathbf{0 9}^{\text {th }}$ SEP, 2021)
\section*{PHYSICS}
\begin{enumerate}

\section*{PHYSICS}

Q1. The equivalent capacitance of the combination shown in the figure is:
$\\$ IMAGE:https://topnow.s3.us-east-1.amazonaws.com/neet/2025_10_11_d2c601cc51a969098abag-01.jpg $\\$

(1) 3 C
(2) 2 C
(3) $C / 2$
(4)$3 c / 2$
Ans. (2)
Sol. $\mathrm{C}_{\mathrm{P}}=\mathrm{C}+\mathrm{C}=2 \mathrm{C}$

Q2. Polar molecules are the molecules:

(1) having zero dipole moment.
(2) acquire a dipole moment only in the presence of electric field due to displacement of charges.
(3) acquire a dipole moment only when magnetic field is absent.
(4) having a permanent electric dipole moment.
Ans. (4)
Sol. Conceptual

Q3. An infinitely long straight conductor carries a current of 5 A as shown. An electron is moving with a speed of $10^{5} \mathrm{~m} / \mathrm{s}$ parallel to the conductor. The perpendicular distance between the electron and the conductor is 20 cm at an instant. Calculate the magnitude of the force experienced by the electron at that instant.
$\\$ IMAGE:https://topnow.s3.us-east-1.amazonaws.com/neet/2025_10_11_d2c601cc51a969098abag-01.jpg $\\$
$\\$ IMAGE:https://topnow.s3.us-east-1.amazonaws.com/neet/2025_10_11_d2c601cc51a969098abag-01.jpg $\\$

(1) $4 \times 10^{-20} \mathrm{~N}$
(2) $8 \pi \times 10^{-20} \mathrm{~N}$
(3) $4 \pi \times 10^{-20} \mathrm{~N}$
(4) $8 \times 10^{-20} \mathrm{~N}$
Ans. (4)
Sol. $\mathrm{F}=\mathrm{BeV}$

$$
=\frac{\mu_{0} \mathrm{i}}{2 \pi \mathrm{r}} \mathrm{ev}
$$

Q4. A thick current carrying cable of radius ' R ' carries current ' I ' uniformly distributed across its cross-section. The variation of magnetic field $B(r)$ due to the cable with the distance ' $r$ ' from the axis of the cable is represented by:

\begin{figure}[h]
\begin{center}
\captionsetup{labelformat=empty}
\caption{1}
(1) $\\$ IMAGE:https://topnow.s3.us-east-1.amazonaws.com/neet/2025_10_11_d2c601cc51a969098abag-02(1).jpg $\\$
\end{center}
\end{figure}

\begin{figure}[h]
\begin{center}
\captionsetup{labelformat=empty}
\caption{2}
(2) $\\$ IMAGE:https://topnow.s3.us-east-1.amazonaws.com/neet/2025_10_11_d2c601cc51a969098abag-02.jpg $\\$
\end{center}
\end{figure}

\begin{figure}[h]
\begin{center}
\captionsetup{labelformat=empty}
\caption{3}
(3) $\\$ IMAGE:https://topnow.s3.us-east-1.amazonaws.com/neet/2025_10_11_d2c601cc51a969098abag-02(4).jpg $\\$
\end{center}
\end{figure}

\begin{figure}[h]
\begin{center}
\captionsetup{labelformat=empty}
\caption{4}
(4) $\\$ IMAGE:https://topnow.s3.us-east-1.amazonaws.com/neet/2025_10_11_d2c601cc51a969098abag-02(3).jpg $\\$
\end{center}
\end{figure}
Ans. (3)
Sol. $r<R \quad B \propto r$
$\mathrm{r}>\mathrm{R}, \mathrm{B} \propto \frac{1}{\mathrm{r}}$

Q5. In a potentiometer circuit a cell of EMF 1.5 V gives balance point at 36 cm length of wire. If another cell of EMF 2.5 V replaces the first cell, then at what length of the wire, the balance point occurs?

(1) 60 cm
(2) 21.6 cm
(3) 64 cm
(4) 62 cm
Ans. (1)
Sol. $\frac{\mathrm{E}_{1}}{\mathrm{E}_{2}}=\frac{\ell_{1}}{\ell_{2}}$

Q6. For a plane electromagnetic wave propagating in x-direction, which one of the following combination gives the correct possible directions for electric field (E) and magnetic field $(B)$ respectively?

(1) $\hat{j}+\hat{k}, \hat{j}+\hat{k}$
(2) $-\hat{j}+\hat{k},-\hat{j}-\hat{k}$
(3) $\hat{j}+\hat{k},-\hat{j}-\hat{k}$
(4) $-\hat{j}+\hat{k},-\hat{j}+\hat{k}$
Ans. (2)
Sol. $\bar{V}=\vec{E} \times \vec{B}$ \& $\vec{E} \cdot \vec{B}=0$

Q7. An inductor of inductance $L$, a capacitor of capacitance $C$ and a resistor of resistance ' $R$ ' are connected in series to an ac source of potential difference ' $V$ ' volts as shown in figure. Potential difference across $\mathrm{L}, \mathrm{C}$ and R is 40 V , 10 V and 40 V , respectively. The amplitude of current flowing through LCR series circuit is $10 \sqrt{2} \mathrm{~A}$. The impedance of the circuit is
$\\$ IMAGE:https://topnow.s3.us-east-1.amazonaws.com/neet/2025_10_11_d2c601cc51a969098abag-02(2).jpg $\\$

(1) $4 \sqrt{2} \Omega$
(2) $5 / \sqrt{2} \Omega$
(3) $4 \Omega$
(4) $5 \Omega$
Ans. (4)
Sol. $i_{\text {rms }}=\frac{V_{\text {rms }}}{Z} \Rightarrow Z=\frac{\sqrt{(40-10)^{2}+40^{2}}}{10}$

Q8. The number of photons per second on an average emitted by the source of monochromatic light of wavelength 600 nm , when it delivers the power of $3.3 \times 10^{-3}$ watt will be : $\left(\mathrm{h}=6.6 \times 10^{-34} \mathrm{Js}\right)$

(1) $10^{18}$
(2) $10^{17}$
(3) $10^{16}$
(4) $10^{15}$
Ans. (3)
Sol. $P=\frac{N}{t} \frac{h c}{\lambda}$

Q9. An electromagnetic wave of wavelength ' $\lambda$ ' is incident on a photosensitive surface of negligible work function. If ' $m$ ' mass is of photoelectron emitted from the surface has deBroglie wavelength $\lambda_{\mathrm{d}}$, then:

(1) $\lambda=\left(\frac{2 \mathrm{~m}}{\mathrm{hc}}\right) \lambda_{\mathrm{d}}{ }^{2}$
(2) $\lambda_{\mathrm{d}}=\left(\frac{2 \mathrm{mc}}{\mathrm{h}}\right) \lambda^{2}$
(3) $\lambda=\left(\frac{2 \mathrm{mc}}{\mathrm{h}}\right) \lambda_{\mathrm{d}}{ }^{2}$
(4) $\lambda=\left(\frac{2 \mathrm{~h}}{\mathrm{mc}}\right) \lambda_{\mathrm{d}}{ }^{2}$
Ans. (3)
Sol. $\mathrm{KE}=\mathrm{E}$

$$
\frac{\mathrm{P}^{2}}{2 \mathrm{~m}}=\frac{\mathrm{hc}}{\lambda}, \frac{\mathrm{~h}^{2}}{\lambda_{\mathrm{d}}{ }^{2} 2 \mathrm{~m}}=\frac{\mathrm{hc}}{\lambda}
$$

Column-I gives certain physical terms


q10. 
{Column-I}
(A) Drift Velocity
(B) Electrical Resistivity
(C) Relaxation Period
(D) Current Density associated with flow of current through a metallic conductor.
Column-II gives some mathematical relations involving electrical quantities. Match ColumnI and Column-II with appropriate relations.

{Column-II}
(P) $\frac{m}{n e^{2} \rho}$
(Q) $n e v_{d}$
(R) $\frac{\mathrm{eE}}{\mathrm{m}} \tau$
(S) $\frac{\mathrm{E}}{\mathrm{J}}$

(1) $(\mathrm{A})-(\mathrm{R}),(\mathrm{B})-(\mathrm{S}),(\mathrm{C})-(\mathrm{P}),(\mathrm{D})-(\mathrm{Q})$
(2) $(\mathrm{A})-(\mathrm{R}),(\mathrm{B})-(\mathrm{S}),(\mathrm{C})-(\mathrm{Q}),(\mathrm{D})-(\mathrm{P})$
(3) $(\mathrm{A})-(\mathrm{R}),(\mathrm{B})-(\mathrm{P}),(\mathrm{C})-(\mathrm{S}),(\mathrm{D})-(\mathrm{Q})$
(4) $(\mathrm{A})-(\mathrm{R}),(\mathrm{B})-(\mathrm{Q}),(\mathrm{C})-(\mathrm{S}),(\mathrm{D})-(\mathrm{P})$
Ans. (1)
Sol. Conceptual

Q11. The escape velocity from the Earth's surface is $v$. The escape velocity from the surface of another planet having a radius, four times that of Earth and same mass density is:

(1) $v$
(2) $2 v$
(3) $3 v$
(4) $4 v$
Ans. (4)
Sol. $\mathrm{V}=\sqrt{2 \mathrm{gR}}=\mathrm{R} \sqrt{\frac{8}{3} \pi \rho \mathrm{G}}, \mathrm{V} \propto \mathrm{R}$

Q12. The velocity of a small ball of mass M and density $d$, when dropped in a container filled with glycerine becomes constant after some time. If the density of glycerine is $\frac{d}{2}$, then the viscous force acting on the ball will be:

(1) $\frac{M g}{2}$
(2) Mg
(3) $\frac{3}{2} \mathrm{Mg}$
(4) 2 Mg
Ans. (1)
Sol. $\mathrm{F}_{\mathrm{V}}=\mathrm{Mg}-\mathrm{F}_{\mathrm{B}}=\mathrm{Mg}-\mathrm{V} \frac{\mathrm{d}}{2} \mathrm{~g}=\mathrm{Mg}-\frac{\mathrm{Mg}}{2}=\frac{\mathrm{Mg}}{2}$

Q13. A body is executing simple harmonic motion with frequency ' $n$ ', the frequency of its potential energy is:

(1) $n$
(2) $2 n$
(3) $3 n$
(4) $4 n$
Ans. (2)
Sol. Conceptual

Q14. Water falls from a height of 60 m at the rate of $15 \mathrm{~kg} / \mathrm{s}$ to operate a turbine. The lasses due to frictional force are $10 \%$ of the input energy. How much power is generated by the turbine? $\left(\mathrm{g}=10 \mathrm{~m} / \mathrm{s}^{2}\right)$

(1) 10.2 kW
(2) 8.1 kW
(3) 12.3 kW
(4) 7.0 kW
Ans. (2)
Sol. $\mathrm{P}=\frac{90}{100} \frac{\mathrm{Mgh}}{\mathrm{t}}$

Q15. A dipole is placed in an electric field as shown. In which direction will it move?
$\\$ IMAGE:https://topnow.s3.us-east-1.amazonaws.com/neet/2025_10_11_d2c601cc51a969098abag-03.jpg $\\$

(1) towards the left as its potential energy will increase.
(2) towards the right as its potential energy will decrease.
(3) towards the left as its potential energy will decrease.
(4) towards the right as its potential energy will increase.
Ans. (2)
Sol. $\mathrm{U}=-\mathrm{PE} \cos \theta,\left(\theta=180^{\circ}\right)$

Q16. A capacitor of capacitance ' C ', is connected across an ac source of voltage $V$, given by $\mathrm{V}=\mathrm{V}_{\mathrm{O}} \sin \omega \mathrm{t}$
The displacement current between he plates of the capacitor, would then be given by:

(1) $I_{d}=V_{0} \omega C \cos \omega t$
(2) $I_{d}=\frac{V_{O}}{\omega C} \cos \omega t$
(3) $I_{d}=\frac{V_{O}}{\omega C} \sin \omega t$
(4) $I_{d}=V_{0} \omega C \sin \omega t$
Ans. (1)
Sol. $I_{d}=C \frac{d v}{d t}=C \frac{d\left(V_{0} \sin \omega t\right)}{d t}=V_{0} \omega C \cos \omega t$

Q17. A cup of coffee cools from $90^{\circ} \mathrm{C}$ to $80^{\circ} \mathrm{C}$ in t minutes, when the room temperature is $20^{\circ} \mathrm{C}$ . The time taken by a similar cup of coffee to cool from $80^{\circ} \mathrm{C}$ to $60^{\circ} \mathrm{C}$ at a room temperature same at $20^{\circ} \mathrm{C}$ is:

(1) $\frac{13}{10} \mathrm{t}$
(2) $\frac{13}{5} t$
(3) $\frac{10}{13} \mathrm{t}$
(4) $\frac{5}{13} \mathrm{t}$
Ans. (2)
Sol. $\frac{\theta_{1}-\theta_{2}}{\mathrm{t}}=-\mathrm{k}\left(\frac{\theta_{1}+\theta_{2}}{2}-\theta_{0}\right)$

Q18. The effective resistance of parallel connection that consists of four wires of equal length, equal area of cross-section and same material is $0.25 \Omega$. What will be the effective resistance if they are connected in series?

(1) $0.25 \Omega$
(2) $0.5 \Omega$
(3) $1 \Omega$
(4) $4 \Omega$
Ans. (4)
Sol. $R_{P}=\frac{R}{n} \Rightarrow 0.25=\frac{R}{4} \Rightarrow R=1 \Omega$
$\mathrm{R}_{\mathrm{S}}=\mathrm{nR}=4 \times 1=4 \Omega$

Q19. Match Column-I and Column-II and choose the correct match from the given choices.

\begin{center}
\begin{tabular}{|l|l|l|l|}
\hline
 & Column-I &  & Column-II 
\hline
(A) & Root mean square speed of gas molecules & (P) & $\frac{1}{3} \mathrm{~nm} \overline{\mathrm{v}}^{2}$ 
\hline
(B) & Pressure exerted by ideal gas & (Q) & $\sqrt{\frac{3 \mathrm{RT}}{\mathrm{M}}}$ 
\hline
(C) & Average kinetic energy of a molecule & (R) & $\frac{5}{2} \mathrm{RT}$ 
\hline
(D) & Total internal energy of 1 mole of a diatomic gas & (S) & $\frac{3}{2} \mathrm{k}_{\mathrm{B}} \mathrm{T}$ 
\hline
\end{tabular}
\end{center}
(1) $(\mathrm{A})-(\mathrm{R}),(\mathrm{B})-(\mathrm{P}),(\mathrm{C})-(\mathrm{S}),(\mathrm{D})-(\mathrm{Q})$
(2) $(\mathrm{A})-(\mathrm{Q}),(\mathrm{B})-(\mathrm{R}),(\mathrm{C})-(\mathrm{S}),(\mathrm{D})-(\mathrm{P})$
(3) $(\mathrm{A})-(\mathrm{Q}),(\mathrm{B})-(\mathrm{P}),(\mathrm{C})-(\mathrm{S}),(\mathrm{D})-(\mathrm{R})$
(4) $(\mathrm{A})-(\mathrm{R}),(\mathrm{B})-(\mathrm{Q}),(\mathrm{C})-(\mathrm{P}),(\mathrm{D})-(\mathrm{S})$
Ans. (3)
Sol. Conceptual

Q20. A small block slides down on a smooth inclined plane, starting form rest at time $t=0$. Let $S_{n}$ be the distance travelled by the block in the interval $t=n-1$ to $t=n$. Then, the ratio $\frac{S_{n}}{S_{n+1}}$ is:

(1) $\frac{2 n-1}{2 n}$
(2) $\frac{2 n-1}{2 n+1}$
(3) $\frac{2 n+1}{2 n-1}$
(4) $\frac{2 n}{2 n-1}$
Ans. (2)
Sol. $\frac{S_{n}}{S_{n+1}}=\frac{g\left(n-\frac{1}{2}\right)}{g\left(n+1-\frac{1}{2}\right)}$

Q21. A radioactive nucleus ${ }_{Z}^{A} \mathrm{X}$ undergoes spontaneous decay in the sequence ${ }_{Z}^{A} X \rightarrow_{Z-1} B \rightarrow_{Z-3} C \rightarrow_{Z-2} D$, where $Z$ is the atomic number of element $X$. The possible decay particles in the sequence are:

(1) $\alpha, \beta^{-}, \beta^{+}$
(2) $\alpha, \beta^{+}, \beta^{-}$
(3) $\beta^{+}, \alpha, \beta$
(4) $\beta^{-}, \alpha, \beta^{+}$
Ans. (3)
Sol. Conceptual

Q22. A screw gauge gives the following readings when used to measure the diameter of a wire Main scale reading : 0 mm
Circular scale reading : 52 divisions
Given that 1 mm on main scale correspond to 100 divisions on the circular scale. The diameter of the wire from the above date is:

(1) 0.52 cm
(2) 0.026 cm
(3) 0.26 cm
(4) 0.052 cm
Ans. (4)
Sol. T.R. $=$ M.S.R $+($ H.S.R $\times$ L.C $)$

Q23. A parallel place capacitor has a uniform electric field ' $\vec{E}$ ' in the space between the plates. If the distance between the plates is 'd' and the area of each plate is 'A', the energy stored in the capacitor is : ( $\varepsilon_{0}=$ permittivity of free space)

(1) $\frac{1}{2} \varepsilon_{0} \mathrm{E}^{2}$
(2) $\varepsilon_{0}$ EAd
(3) $\frac{1}{2} \varepsilon_{0} \mathrm{E}^{2} \mathrm{Ad}$
(4) $\frac{\mathrm{E}^{2} \mathrm{Ad}}{\varepsilon_{0}}$
Ans. (3)
Sol. $\mathrm{U}=\frac{1}{2} \mathrm{Cv}^{2}=\frac{1}{2} \varepsilon_{0} \mathrm{E}^{2} \mathrm{Ad}$

Q24. The electron concentration in an n-type semiconductor is the same as hole concentration in a p-type semiconductor. An external field (electric) is applied across each of them. Compare the current in them.

(1) current in n-type $=$ current in p-type.
(2) current in p-type > current in n-type.
(3) current in n-type > current in p-type.
(4) No current will flow in p-type, current will only flow in n-type.
Ans. (3)
Sol. Mobility of electrons is more than holes

Q25. A convex lens 'A' of focal length 20 cm and a concave lens ' B ' of focal length 5 cm are kept along the same axis with a distance 'd' between them. If a parallel beam of light falling on 'A' leaves ' B ' as a parallel beam, then the distance ' $d$ ' in cm will be:

(1) 25
(2) 15
(3) 50
(4) 30
Ans. (2)_
Sol. $\quad d=f_{1}-f_{2}$

Q26. Find the value of the angle of emergence from the prism. Refractive index of the glass is $\sqrt{3}$.

(1) $60^{\circ}$
(2) $30^{\circ}$
(3) $45^{\circ}$
(4) $90^{\circ}$
Ans.(1)
Sol. $\quad \mathbf{i}_{\mathbf{i}}=\mathbf{r}_{1}=0$

$$
r_{2}=30^{\circ}, \mu=\frac{\sin i_{2}}{\sin r_{2}}
$$

q27.  Consider the following statements (A) and (B) and identify the correct answer.
(A) A Zener diode is connected in reverse bias, when used as a voltage regulator.
(B) The potential barrier of $\mathrm{p}-\mathrm{n}$ junction lies between 0.1 V to 0.3 V .

(1) (A) and (B) both are correct
(2) (A) and (B) both are incorrect
(3) (A) is correct and (B) is incorrect
(4) (A) is incorrect but (B) is correct
Ans. (3)
Sol. Conceptual

Q28. Two charged spherical conductors of radius $R_{1}$ and $R_{2}$ are connect by a wire. Then the ratio of surface charge densities of the spheres ( $\sigma_{1} / \sigma_{2}$ ) is:

(1) $\frac{R_{1}}{R_{2}}$
(2) $\frac{R_{2}}{R_{1}}$
(3) $\sqrt{\left(\frac{\mathrm{R}_{1}}{\mathrm{R}_{2}}\right)}$
(4) $\frac{R_{1}^{2}}{R_{2}^{2}}$
Ans. (2)
Sol. $\frac{\sigma_{1}}{\sigma_{2}}=\frac{\mathbf{q}_{1}}{\mathbf{q}_{2}}\left(\frac{\mathbf{R}_{2}}{\mathbf{R}_{1}}\right)^{2}$

$$
\mathrm{k} \frac{\mathrm{q}_{1}}{\mathrm{R}_{1}}=\mathrm{v}=\mathrm{k} \frac{\mathrm{q}_{2}}{\mathrm{R}_{2}} \Rightarrow \frac{\mathrm{q}_{1}}{\mathrm{q}_{2}}=\frac{\mathrm{R}_{1}}{\mathrm{R}_{2}}
$$

Q29. If force $[\mathrm{F}]$, acceleration $[\mathrm{A}]$ and time $[\mathrm{T}]$ are chosen as the fundamental physical quantities. Find the dimensions of energy.

(1) $[\mathrm{F}][\mathrm{A}][\mathrm{T}]$
(2) $[\mathrm{F}][\mathrm{A}]\left[\mathrm{T}^{2}\right]$
(3) $[\mathrm{F}][\mathrm{A}]\left[\mathrm{T}^{-1}\right]$
(4) $[\mathrm{F}]\left[\mathrm{A}^{-1}\right][\mathrm{T}]$
Ans. (2)
Sol. Conceptual

Q30. If E and G respectively denote energy and gravitational constant, then $\frac{E}{G}$ has the dimensions of:

(1) $\left[\mathrm{M}^{2}\right]\left[\mathrm{L}^{-1}\right]\left[\mathrm{T}^{0}\right]$
(2) $[\mathrm{M}]\left[\mathrm{L}^{-1}\right]\left[\mathrm{T}^{-1}\right]$
(3) $[\mathrm{M}]\left[\mathrm{L}^{0}\right]\left[\mathrm{T}^{0}\right]$
(4) $\left[\mathrm{M}^{2}\right]\left[\mathrm{L}^{-2}\right]\left[\mathrm{T}^{-1}\right]$
Ans. (1)

Sol. Conceptual

Q31. A lens of large focal length and large aperture is best suited as an objective of an astronomical telescope since

(1) A larger aperture contributes to the equality and visibility of the images
(2) A large area of the objective ensures better light gathering power
(3) A large aperture provides a better resolution
(4) All of the above
Ans. (4)
Sol. Conceptual

Q32. A nucleus with mass number 240 breaks into two fragments each of mass number 120, the binding energy per nucleon of unfragmented nuclei is 7.6 MeV while that of fragments is 8.5 MeV . The total gain the Binding Energy in the process is

(1) 0.9 MeV
(2) 9.4 MeV
(3) 804 MeV
(4)216 MeV
Ans.(4)
Sol. B.E $=(2 \times 120 \times 8.5)-(240 \times 7.6) \mathrm{MeV}$

Q33. A particle is released from height S from the surface to the Earth. At a certain height its kinetic energy is three times its potential energy. The height from the surface of earth and the speed of the particle at that instant are respectively

(1) $\frac{\mathrm{S}}{4}, \frac{3 \mathrm{gS}}{2}$
(2) $\frac{\mathrm{S}}{4}, \frac{\sqrt{3 \mathrm{gS}}}{2}$
(3) $\frac{\mathrm{S}}{2}, \frac{\sqrt{3 \mathrm{gS}}}{2}$
(4) $\frac{\mathrm{S}}{4}, \sqrt{\frac{3 \mathrm{gS}}{2}}$
Ans. (4)
Sol. $\frac{\mathrm{PE}}{\mathrm{KE}}=\frac{\mathrm{h}}{\mathrm{S}-\mathrm{h}} \Rightarrow \frac{\mathrm{PE}}{3 \mathrm{PE}}=\frac{\mathrm{h}}{\mathrm{S}-\mathrm{h}} \Rightarrow \mathrm{h}=\frac{\mathrm{S}}{4}$
$V=\sqrt{2 g \frac{3 s}{4}}$

Q34. The half-life of a radioactive nuclide is 100 hours. The fraction of original activity that will remain after 150 hours would be

(1) $1 / 2$
(2) $\frac{1}{2 \sqrt{2}}$
(3) $\frac{2}{3}$
(4) $\frac{2}{3 \sqrt{2}}$
Ans. (2)
Sol. $\mathrm{N}=\mathrm{N}_{\mathrm{O}}\left(\frac{1}{2^{\mathrm{t} / \mathrm{T}}}\right), \mathrm{t}=\frac{3 \mathrm{~T}}{2}$

Q35. A spring is stretched by 5 cm by a force 10 N . The time period of the oscillations when a mass of 2 kg is suspended by it is

(1) 0.0628 s
(2) 6.28 s
(3) 3.14 s
(4) 0.628 s
Ans. (4)
Sol. $\mathrm{K}=\frac{\mathrm{F}}{\mathrm{x}}, \mathrm{T}=2 \pi \sqrt{\frac{\mathrm{~m}}{\mathrm{k}}}$


Q36. A series LCR circuit containing 5.0 H inductor, $80 \mu \mathrm{~F}$ capacitor and $40 \Omega$ resistor is connected to 230 V variable frequency ac source. The angular frequencies of the at which power transferred to the circuit is half the power at the resonant angular frequency are likely to be

(1) $25 \mathrm{rad} / \mathrm{s}$ and $75 \mathrm{rad} / \mathrm{s}$
(2) $50 \mathrm{rad} / \mathrm{s}$ and $25 \mathrm{rad} / \mathrm{s}$
(3) $46 \mathrm{rad} / \mathrm{s}$ and $54 \mathrm{rad} / \mathrm{s}$
(4) $42 \mathrm{rad} / \mathrm{s}$ and $58 \mathrm{rad} / \mathrm{s}$
Ans. (3)
Sol. $\omega=\left(\sqrt{\left(\frac{\mathrm{R}}{2 \mathrm{~L}}\right)^{2}+\frac{1}{\mathrm{LC}}}\right) \pm \frac{\mathrm{R}}{2 \mathrm{~L}}$

Q37. The conducting circular loops of radii $\mathrm{R}_{1}$ and $R_{2}$ are placed in the same plane with their centres coinciding. If $R_{1} \gg R_{2}$, the mutual inductance $M$ between them will be directly proportional to

(1) $\frac{R_{1}}{R_{2}}$
(2) $\frac{R_{2}}{R_{1}}$
(3) $\frac{R_{1}^{2}}{R_{2}}$
(4) $\frac{R_{2}^{2}}{R_{1}}$
Ans. (4)
Sol. $\mathrm{N} \phi=\mathrm{Mi} \Rightarrow \mathrm{NBA}_{2}=\mathrm{Mi}$

$$
\frac{\mathrm{N} \mu_{0} \mathrm{i}}{2 \mathrm{R}_{1}} \pi \mathrm{R}^{2}=\mathrm{Mi}
$$

Q38. A ball of mass 0.15 kg is dropped from a height 10 m , strikes the ground and rebounds to the same height. The magnitude of impulse imparted to the ball is $\left(g=10 \mathrm{~m} / \mathrm{s}^{2}\right)$ nearly

(1) $0 \mathrm{~kg} \mathrm{~m} / \mathrm{s}$
(2) $4.2 \mathrm{~kg} \mathrm{~m} / \mathrm{s}$
(3) $2.1 \mathrm{~kg} \mathrm{~m} / \mathrm{s}$
(4) $1.4 \mathrm{~kg} \mathrm{~m} / \mathrm{s}$
Ans. (2)
Sol. $\mathrm{J}=\Delta \mathrm{P}=\mathrm{m}\left(\sqrt{2 \mathrm{gh}_{2}}+\sqrt{2 \mathrm{gh}_{1}}\right),\left(\mathrm{h}_{1}=\mathrm{h}_{2}\right)$

Q39. A point object is placed at a distance of 60 cm from a convex lens of focal length 30 cm . If a plane mirror were put perpendicular to the principal axis of the lens and at a distance of 40 cm from it, the final image would be formed at a distance of
$\\$ IMAGE:https://topnow.s3.us-east-1.amazonaws.com/neet/2025_10_11_d2c601cc51a969098abag-07.jpg $\\$

(1) 20 cm from the lens, it would be a real image
(2) 30 cm from the lens, it would be a real image
(3) 30 cm from the plane mirror, it would be virtual image
(4) 20 cm from the plane mirror, it would be a virtual image
Ans. (1)
Sol. $\frac{1}{V}-\frac{1}{u}=\frac{1}{f} \Rightarrow \frac{1}{V}-\frac{1}{-60}=\frac{1}{30} \Rightarrow V=60 \mathrm{~cm}$
$\\$ IMAGE:https://topnow.s3.us-east-1.amazonaws.com/neet/2025_10_11_d2c601cc51a969098abag-07(2).jpg $\\$


This virtual image acts as an object for plane mirror and final image is formed at 20 cm from the lens, it would be a real image

Q40. A uniform conducting wire of length 12 a and resistance ' $R$ ' is wound up as a current carrying coil in the shape of
(i) An equilateral triangle of side 'a'
(ii) A square of side 'a'

The magnetic dipole moments of the coil in each case respectively are

(1) $\sqrt{3} I a^{2}$ and $3 I a^{2}$
(2) $3 I a^{2}$ and $I a^{2}$
(3) $3 I a^{2}$ and $4 I a^{2}$
(4) $4 I a^{2}$ and $3 I a^{2}$
Ans. (1)
Sol. $\mathrm{M}=\mathrm{niA}$
For equilateral triangle $\mathrm{n}=4, \mathrm{~A}=\frac{1}{2} \mathrm{a} \frac{\sqrt{3}}{2} \mathrm{a}$
For square $\mathrm{n}=3, \mathrm{~A}=\mathrm{a}^{2}$

Q41. A car starts from rest and accelerates at 5 $\mathrm{m} / \mathrm{s}^{2}$. At $\mathrm{t}=4 \mathrm{~s}$, a ball is dropped out of a window by a person sitting in the car. What is the velocity and acceleration of the ball at $\mathrm{t}=6 \mathrm{~s}$ ? (Take $\mathrm{g}=10 \mathrm{~m} / \mathrm{s}^{2}$ )

(1) $20 \mathrm{~m} / \mathrm{s}, 5 \mathrm{~m} / \mathrm{s}^{2}$
(2)$20 \mathrm{~m} / \mathrm{s}, 0$
(3) $20 \sqrt{2} \mathrm{~m} / \mathrm{s}, 0$
(4) $20 \sqrt{2} \mathrm{~m} / \mathrm{s}, 10 \mathrm{~m} / \mathrm{s}^{2}$
Ans. (4)
Sol. $\mathrm{V}_{\mathrm{x}}=\mathrm{a}_{\mathrm{x}} \mathrm{t}=20$

$$
\begin{aligned}
& V_{y}=g t=10(6-4) 
& V=\sqrt{V_{x}^{2}+V_{y}^{2}}, a=g
\end{aligned}
$$

Q42. Three resistor having resistances $\mathrm{r}_{1}, \mathrm{r}_{2}$ and $\mathrm{r}_{3}$ are connected as shown in the given circuit. The ratio $\frac{i_{3}}{i_{1}}$ of currents in terms of resistance use in the circuit is
$\\$ IMAGE:https://topnow.s3.us-east-1.amazonaws.com/neet/2025_10_11_d2c601cc51a969098abag-07(1).jpg $\\$

(1) $\frac{r_{1}}{r_{2}+r_{3}}$
(2) $\frac{r_{2}}{r_{2}+r_{3}}$
(3) $\frac{r_{1}}{r_{1}+r_{2}}$
(4) $\frac{r_{2}}{r_{1}+r_{3}}$
Ans. (2)
Sol. $i_{3}=i_{1}\left(\frac{r_{2}}{r_{2}+r_{3}}\right)$

Q43. A particle of mass ' $m$ ' is projected with a velocity $\mathrm{v}=\mathrm{kV}_{\mathrm{e}}(\mathrm{k}<1)$ from the surface of earth ( $\mathrm{V}_{\mathrm{e}}=$ escape velocity)
The maximum height above the surface reached by the particle is

(1) $\mathrm{R}\left(\frac{\mathrm{k}}{1-\mathrm{k}}\right)^{2}$
(2) $\mathrm{R}\left(\frac{\mathrm{k}}{1+\mathrm{k}}\right)^{2}$
(3)  $\frac{R^{2} k}{1+k}$
(4) $\frac{\mathrm{Rk}^{2}}{1-\mathrm{k}^{2}}$
Ans. (4)
Sol. $\frac{1}{2} \mathrm{mK}^{2} \mathrm{~V}_{\mathrm{e}}^{2}=\frac{\mathrm{mgh}}{\left(1+\frac{\mathrm{h}}{\mathrm{R}}\right)}$

Q44. A step down transformer connected to an ac mains supply of 220 V is made to operated at $11 \mathrm{~V}, 44 \mathrm{~W}$ lamp. Ignoring power losses in the transformer, what is the current in the primary circuit.

(1) 0.2 A
(2) 0.4 A
(3) 2 A
(4) 4 A
Ans. (1)
Sol. $\mathrm{P}_{\text {Out }}=\mathrm{P}_{\mathrm{In}} \Rightarrow 44=\mathbf{i}_{\mathrm{p}}(220)$

Q45. Twenty seven drops of same size are charge at 220 V each. They combine to form a bigger drop. Calculate the potential of the bigger drop

(1) 660 V
(2) 1320 V
(3) 1520 V
(4) 1980 V
Ans. (4)
Sol. $\mathrm{V}^{\prime}=\mathrm{n}^{2 / 3} \mathrm{~V}$

Q46. For the given circuit, the input digital signals are applied at the terminals $\mathrm{A}, \mathrm{B}$ and C . What would be the out put at the terminaly?
$\\$ IMAGE:https://topnow.s3.us-east-1.amazonaws.com/neet/2025_10_11_d2c601cc51a969098abag-08(3).jpg $\\$
(1) $\\$ IMAGE:https://topnow.s3.us-east-1.amazonaws.com/neet/2025_10_11_d2c601cc51a969098abag-08(4).jpg $\\$

(2) $\\$ IMAGE:https://topnow.s3.us-east-1.amazonaws.com/neet/2025_10_11_d2c601cc51a969098abag-08(1).jpg $\\$

(3) $\\$IMAGE:https://topnow.s3.us-east-1.amazonaws.com/neet/2025_10_11_d2c601cc51a969098abag-08(2).jpg $\\$

(4) $\\$ IMAGE:https://topnow.s3.us-east-1.amazonaws.com/neet/2025_10_11_d2c601cc51a969098abag-08.jpg $\\$
Ans. (3)
Sol. $Y=A \cdot B+\overline{B \cdot C}$

Q47. A particle moving in a circle of radius $R$ with a uniform speed takes a time T to complete one revolution.
If this particle were projected with the same speed at an angle ' $\theta$ ' to the horizontal, the maximum height attained by it equals 4 R . The angle of projection, $\theta$, is then given by

(1) $\theta=\boldsymbol{\operatorname { c o s }}^{-1}\left(\frac{\mathbf{g T}^{2}}{\pi^{2} \mathrm{R}}\right)^{1 / 2}$
(3) $\theta=\boldsymbol{\operatorname { c o s }}^{-1}\left(\frac{\pi^{2} \mathrm{R}}{\mathbf{g T}^{2}}\right)^{1 / 2}$
(3) $\theta=\sin ^{-1}\left(\frac{\pi^{2} \mathrm{R}}{g \mathrm{~T}^{2}}\right)^{1 / 2}$
(4) $\theta=\sin ^{-1}\left(\frac{2 g T^{2}}{\pi^{2} R}\right)^{1 / 2}$
Ans. (4)
Sol. $u=\frac{2 \pi R}{T}, H=\frac{u^{2} \sin ^{2} \theta}{2 g}$

Q48. From a circular ring of mass ' M ' and radius ' R ' an arc corresponding to a $90^{\circ}$ sector is removed. The moment of inertia of the remaining part of the ring about an axis passing through the centre of the ring and perpendicular to the plane of the rings ' K ' time ' $\mathrm{MR}^{2}$ '. Then the value of ' K ' is

(1) $\frac{3}{4}$
(2) $\frac{7}{8}$
(3) $\frac{1}{4}$
(4) $\frac{1}{8}$
Ans. (1)
Sol. $\mathrm{I}_{\text {Remaining }}=\mathrm{I}_{\text {original }}-\mathrm{I}_{\text {removed }} \Rightarrow \mathrm{MR}^{2}-\frac{\mathrm{MR}^{2}}{4}$

Q49. A uniform rod of length 200 cm and mass 500 g is balanced on the wedge placed at 40 cm mark. A mass of 20 kg is suspended from the rod at 20 cm and another unknown mass ' m ' is suspended from the rod at 160 cm mark as shown in the figure. Find the value of ' $m$ ' such that the rod is in equilibrium. ( $g=10 \mathrm{~m} / \mathrm{s}^{2}$ )
$\\$ IMAGE:https://topnow.s3.us-east-1.amazonaws.com/neet/2025_10_11_d2c601cc51a969098abag-08(5).jpg $\\$

(1) $\frac{1}{2} \mathrm{~kg}$
(2) $\frac{1}{3} \mathrm{~kg}$
(3) $\frac{1}{6} \mathrm{~kg}$
(4) $\frac{1}{12} \mathrm{~kg}$
Ans. (4)
Sol. Net torque is zero

$$
2 g(20)=m g(120)+(0.5 g)(60)
$$

Q50.
 In the product
$\vec{F}=q(\vec{v} \times \vec{B})=q \vec{v} \times\left(B \hat{i}+B \hat{j}+B_{0} \hat{k}\right)$
For $q=1$ and $\vec{v}=2 \hat{i}+4 \hat{j}+6 \hat{k}$ and
$\overrightarrow{\mathrm{F}}=4 \hat{\mathrm{i}}-20 \hat{\mathrm{j}}+12 \hat{\mathrm{k}}$
What will be the complete expression for $\overrightarrow{\mathrm{B}}$ ?

(1) $-8 \hat{\mathrm{i}}-8 \hat{\mathrm{j}}-6 \hat{\mathrm{k}}$
(2) $-6 \hat{i}-6 \hat{j}-8 \hat{k}$
(3) $8 \hat{i}+8 \hat{j}-6 \hat{k}$
(4) $6 \hat{i}+6 \hat{j}-8 \hat{k}$
Ans. (2)
Sol. $\vec{F}=q(\vec{v} \times \vec{B})$

\section*{CHEMISTRY}

Q51. The correct sequence of bond enthalpy of $\mathrm{C}-\mathrm{X}$ bond is:

(1) $\mathrm{CH}_{3}-\mathrm{F}<\mathrm{CH}_{3}-\mathrm{Cl}<\mathrm{CH}_{3}-\mathrm{Br}<\mathrm{CH}_{3}-\mathrm{I}$
(2) $\mathrm{CH}_{3}-\mathrm{F}>\mathrm{CH}_{3}-\mathrm{Cl}>\mathrm{CH}_{3}-\mathrm{Br}>\mathrm{CH}_{3}-\mathrm{I}$
(3) $\mathrm{CH}_{3}-\mathrm{F}<\mathrm{CH}_{3}-\mathrm{Cl}>\mathrm{CH}_{3}-\mathrm{Br}>\mathrm{CH}_{3}-\mathrm{I}$
(4) $\mathrm{CH}_{3}-\mathrm{Cl}>\mathrm{CH}_{3}-\mathrm{F}>\mathrm{CH}_{3}-\mathrm{Br}>\mathrm{CH}_{3}-\mathrm{I}$
Ans. (2)
Sol. Conceptual (Bond energy or enthalpy $\frac{1}{\propto}$ to bond length)

Q52. Which one of the following methods can be used to obtain highly pure metal which is liquid at room temperature?

(1) Electrolysis
(2) Chromatography
(3) Distillation
(4) Zone refining
Ans. (3)
Sol. Conceptual

Q53. The correct option for the number of body centred unit cell in all 14 types of Bravais lattice unit cells is:

(1) 7
(2) 5
(3) 2
(4) 3

Ans. (4)
Sol. Body centred unit cell is possible in (i) cubic (ii) Tetragonal (iii) Orthorhombic

Q54. Among the following alkaline earth metal halides, one which is covalent and soluble in organic solvents is:

(1) Calcium chloride
(2) Strontium chloride
(3) Magnesium chloride
(4) Beryllium chloride
Ans. (4)
Sol. Conceptual

Q55. $\operatorname{Zr}(Z=40)$ and $\operatorname{Hf}(Z=72)$ have similar atomic and ionic radii because of:

(1) belonging to same group
(2) diagonal relationship
(3) lanthanoid contraction
(4) having similar chemical properties
Ans. (3)
Sol. Conceptual

Q56. The maximum temperature than can be achieved in blast furnace is:

(1) upto 1200 K
(2) upto 2200 K
(3) upto 1900 K
(4) upto 5000 K
Ans. (2)
Sol. Conceptual

Q57. What is the IUPAC name of the organic compound formed in the following chemical reaction?
Acetone $\xrightarrow[\text { (ii) } \mathrm{H}_{2} \mathrm{O}, \mathrm{H}^{+}]{\text {(i) } \mathrm{C}_{2} \mathrm{H}_{5} \mathrm{MgBr} \text {,dry Ether }}$ product

(1) 2-methyl propan-2-ol
(2) pentan-2-ol
(3) pentan-3-ol
(4) 2-methyl butan-2-ol
Ans. (4)
Sol. Ketones on reaction with Grignard reagent gives tertiary alcohals.
$\\$ IMAGE:https://topnow.s3.us-east-1.amazonaws.com/neet/smile-2be740d1ffaae920ce5d33d7a00d4b8b188abc89.png $\\$

Q58. Which one of the following polymers is prepared by addition polymerisation?

(1) Teflon
(2) Nylon-66
(3) Novolac
(4) Dacron
Ans. (1)
Sol. Teflon
$\\$ IMAGE:https://topnow.s3.us-east-1.amazonaws.com/neet/2025_10_11_d2c601cc51a969098abag-09.jpg $\\$

Q59. Right option for the number of tetrahedral and octahedral voids in hexagonal primitive unit cell are:

(1) 8,4
(2) 6,12
(3) 2,1
(4) 12,6
Ans. (4)
Sol. Hexagonal primitive unit cell has ' 6 ' atoms per unit cell $Z=6$
$\therefore$ Tetrahedral voids $=2 \mathrm{~N}=12$
Octahedral voids $=\mathrm{N}=6$

Q60. Statement-I : Acid strength increases in the order given as $\mathrm{HF} \ll \mathrm{HCl} \ll \mathrm{HBr} \ll \mathrm{HI}$
Statement - II : As the size of the elements $\mathrm{F}, \mathrm{Cl}, \mathrm{Br}, \mathrm{I}$ increases down the group, the bond strength of $\mathrm{HF}, \mathrm{HCl}, \mathrm{HBr}$ and HI decreases and so the acid strength increases.

(1) Both Statement I and Statement II are true
(2) Both Statement I and Statement II are false
(3) Statement I is correct but Statement II is false
(4) Statement I is incorrect but Statement II is true
Ans. (1)
Sol. Conceptual

Q61. The incorrect Statement among the following is:

(1) Actinoid contraction is greater for element to element than Lanthanoid contraction.
(2) Most of the trivalent Lanthanoid ions are colorless in the solid state.
(3) Lanthanoids are good conductors of heat and electricity.
(4) Actinoids are highly reactive metals especially when finely divided.
Ans. (2)
Sol. Conceptual

Q62. The major product of the following chemical reaction is:
$\\$ IMAGE:https://topnow.s3.us-east-1.amazonaws.com/neet/2025_10_11_d2c601cc51a969098abag-10.jpg $\\$

(1).IMAGE:https://topnow.s3.us-east-1.amazonaws.com/neet/smile-9e281eefaf4628565377ddb8433367543bfc4214.png $\\$
(2) $\\$ IMAGE:https://topnow.s3.us-east-1.amazonaws.com/neet/smile-41cf775cfe543f86d0a77cbae4a6856af390fbfc.png $\\$
(3) $\\$ IMAGE:https://topnow.s3.us-east-1.amazonaws.com/neet/smile-c8e75802f1e0558b91f20ceff6370531c8cfee70.png $\\$
(4) $\\$ IMAGE:https://topnow.s3.us-east-1.amazonaws.com/neet/smile-60d9fe29e0cb00e2f1e7c0feb77cd8f9c896ab53.png $\\$
Ans. (1)
Sol. $\\$ MAGE:https://topnow.s3.us-east-1.amazonaws.com/neet/2025_10_11_d2c601cc51a969098abag-10(1).jpg $\\$

This mechanism is free radical and it follows antimarkonicov's rule.

Q63. The structures of beryllium chloride in solid state and vapour phase, are:

(1) Chain and dimer, respectively
(2) Linear in both
(3) Dimer and Linear, respectively
(4) Chain in both
Ans. (1)
Sol. Conceptual (NCERT text book)

Q64. Given below are two statements:

Statement - I : Aspirin and Paracetamol belong to the class of narcotic analgesics.
Statement - II : Morphine and Heroin are nonnarcotic analgesics.

(1) Both Statement I and Statement II are true
(2) Both Statement I and Statement II are false
(3) Statement I is correct but Statement II is false
(4) Statement I is incorrect but Statement II is true
Ans. (2)
Sol. Conceptual

Q65  An organic compound contain $78 \%$ (by wt.) carbon and remaining percentage of hydrogen. The right option for the empirical formula of this compound is: (Atomic weight of C is 12, H is 1 )

(1) CH
(2) $\mathrm{CH}_{2}$
(3) $\mathrm{CH}_{3}$
(4) $\mathrm{CH}_{4}$
Ans. (3)
Sol. C:H
$\frac{78}{12}: \frac{22}{1}$
$6.5: 22$
$\frac{6.5}{6.5}: \frac{22}{6.5}$
1:3.38
$\cong 1: 3$
$\therefore \mathrm{CH}_{3}$

Q66. The correct structure of 2,6-dimethyl-dec-4ene is:
(1) $\\$ IMAGE:https://topnow.s3.us-east-1.amazonaws.com/neet/smile-1f0d47acf48b4edfbfab61bb25c664da888ea3c5.png $\\$
(2) $\\$ IMAGE:https://topnow.s3.us-east-1.amazonaws.com/neet/smile-8c31619410dbda7a5d7378f4b9e374659bf792d9.png $\\$
(3) $\\$ IMAGE:https://topnow.s3.us-east-1.amazonaws.com/neet/smile-6859ed8af8b1d975991b7f05f59e3c0d9b089e85.png $\\$
(4) $\\$ IMAGE:https://topnow.s3.us-east-1.amazonaws.com/neet/smile-660d850b39647b85ccf9357736bdf95bbedd92d9.png $\\$
Ans. (1)
Sol. Based on IUPAC rules.

Q67. The major product formed in dehydrohalogenation reaction of 2-Bromo pentane is Pent-2-ene. This product formation is based on?

(1) Sayzeffs rule
(2)Hund's rule
(3) Hofmann rule
(4) Huckel's rule
Ans. (1)
Sol. Dehydro halogenation followed by Saytzeff's rule.
(1) $\\$ IMAGE:https://topnow.s3.us-east-1.amazonaws.com/neet/2025_10_11_d2c601cc51a969098abag-11.jpg $\\$
(2) $\\$ IMAGE:https://topnow.s3.us-east-1.amazonaws.com/neet/2025_10_11_d2c601cc51a969098abag-11(3).jpg $\\$
(3) $\\$ IMAGE:https://topnow.s3.us-east-1.amazonaws.com/neet/2025_10_11_d2c601cc51a969098abag-11(2).jpg $\\$
(4) $\\$ IMAGE:https://topnow.s3.us-east-1.amazonaws.com/neet/2025_10_11_d2c601cc51a969098abag-11(1).jpg $\\$

Q68. For a reaction $\mathrm{A} \rightarrow \mathrm{B}$ enthalpy of reaction is $-4.2 \mathrm{kJmol}^{-1}$ and enthalpy of activation is $9.6 \mathrm{kJmol}^{-1}$. The correct potential energy profile for the reaction is shown in option.

(1) $\\$ IMAGE:https://topnow.s3.us-east-1.amazonaws.com/neet/2025_10_11_d2c601cc51a969098abag-11(1).JPG $\\$
(2) $\\$ IMAGE:https://topnow.s3.us-east-1.amazonaws.com/neet/2025_10_11_d2c601cc51a969098abag-11(2).JPG $\\$
(3) $\\$ IMAGE:https://topnow.s3.us-east-1.amazonaws.com/neet/2025_10_11_d2c601cc51a969098abag-11(3).JPG $\\$
(4) $\\$ IMAGE:https://topnow.s3.us-east-1.amazonaws.com/neet/2025_10_11_d2c601cc51a969098abag-11.JPG $\\$
Ans. (2)
Sol. $\Delta \mathrm{H}=-\mathrm{ve}_{\text {(exothermic) }}$
For exothermic reaction $\mathrm{PE}_{(\mathrm{R})}>\mathrm{PE}_{(\mathrm{P})}$
$\left(\mathrm{E}_{\mathrm{a}}\right)_{\mathrm{f}}<\left(\mathrm{E}_{\mathrm{a}}\right)_{\mathrm{b}}$

Q69. Ethylene diaminetetraacetate (EDTA) ion is:

(1) Hexadentate ligand with four ' $O$ ' and two ' $N$ ' donor atoms.
(2) Unidentate ligand
(3) Bidentate ligand with two 'N' donor atoms.
(4) Tridentate ligand with three 'N' donor atoms.
Ans. (1)
Sol. Conceptual

Q70. Noble gases are named because of their inertness towards reactivity. Identify an incorrect statement about them.

(1)Noble gases are sparingly soluble in water
(2) Noble gases have very high melting and boiling points.
(3) Noble gases have weak dispersion forces.
(4) Noble gases have large positive values of electron gain enthalpy.
Ans. (2)
Sol. Conceptual

Q71. 
Which of the following reactions is the metal displacement reaction? Choose the right option.

(1) $2 \mathrm{KClO}_{3} \xrightarrow{\Delta} 2 \mathrm{KCl}+3 \mathrm{O}_{2}$
(2) $\mathrm{Cr}_{2} \mathrm{O}_{3}+2 \mathrm{Al} \xrightarrow{\Delta} \mathrm{Al}_{2} \mathrm{O}_{3}+2 \mathrm{Cr}$
(3) $\mathrm{Fe}+2 \mathrm{HCl} \rightarrow \mathrm{FeCl}_{2}+\mathrm{H}_{2} \uparrow$
(4) $2 \mathrm{~Pb}\left(\mathrm{NO}_{3}\right)_{2} \rightarrow 2 \mathrm{PbO}+4 \mathrm{NO}_{2}+\mathrm{O}_{2} \uparrow$
Ans. (2)
Sol. Conceptual

Q72. The compound which shows metamerism is:

(1) $\mathrm{C}_{5} \mathrm{H}_{12}$
(2) $\mathrm{C}_{3} \mathrm{H}_{8} \mathrm{O}$
(3) $\mathrm{C}_{3} \mathrm{H}_{6} \mathrm{O}$
(4) $\mathrm{C}_{4} \mathrm{H}_{10} \mathrm{O}$
Ans. (4)
Sol. $\mathrm{C}_{4} \mathrm{H}_{10} \mathrm{O}$ - shows metamerism.

$$
\begin{aligned}
& \mathrm{CH}_{3}-\mathrm{O}-\mathrm{C}_{3} \mathrm{H}_{7} 
& \mathrm{C}_{2} \mathrm{H}_{5}-\mathrm{O}-\mathrm{C}_{2} \mathrm{H}_{5}
\end{aligned}
$$

Isomers with same molecular formula and same functional group but which shows difference in attached alkyl groups to the divalent atom or group.

Q73. The RBC deficiency is deficiency disease of

(1) Vitamin $B_{12}$
(2) Vitamin $B_{6}$
(3) Vitamin $B_{1}$
(4) Vitamin $B_{2}$
Ans. (1)
Sol. RBC deficiency disease is vitamin $\mathrm{B}_{12}$

Q74. Dihedral angle of least stable conformer of ethane is:

(1) $120^{\circ}$
(2) $180^{\circ}$
(3) $60^{\circ}$
(4) $0^{\circ}$
Ans. (4)
Sol. Least stable conformer is eclipsed with dihedral angle is zero

Q75. Tritium a radioactive isotope of hydrogen, emits which of the following particles?

(1) Beta ( $\beta^{-}$)
(2) Alpha ( $\alpha$ )
(3) Gamma ( $\gamma$ )
(4) Neutron (n)
Ans. (1)
Sol. Conceptual

Q76. Choose the correct option for graphical representation of Boyle's law, which shows graph of pressure vs volume of a gas at different temperature:

\begin{figure}[h]
\begin{center}
\captionsetup{labelformat=empty}
\caption
(1) $\\$ IMAGE:https://topnow.s3.us-east-1.amazonaws.com/neet/2025_10_11_d2c601cc51a969098abag-12(1).jpg $\\$
\end{center}
\end{figure}

\begin{figure}[h]
\begin{center}
\captionsetup{labelformat=empty}
\caption
(2) $\\$ IMAGE:https://topnow.s3.us-east-1.amazonaws.com/neet/2025_10_11_d2c601cc51a969098abag-12(2).jpg $\\$
\end{center}
\end{figure}

\begin{figure}[h]
\begin{center}
\captionsetup{labelformat=empty}
\caption
(3) $\\$ IMAGE:https://topnow.s3.us-east-1.amazonaws.com/neet/2025_10_11_d2c601cc51a969098abag-12.jpg $\\$
\end{center}
\end{figure}

\begin{figure}[h]
\begin{center}
\captionsetup{labelformat=empty}
\caption
(4) $\\$ IMAGE:https://topnow.s3.us-east-1.amazonaws.com/neet/2025_10_11_d2c601cc51a969098abag-12(3).jpg $\\$
\end{center}
\end{figure}
Ans. (4)
Sol. $\quad \mathrm{P} \propto \mathrm{T}$
P $\propto \frac{1}{\text { V }}$

Q77. The molar conductance of $\mathrm{NaCl}, \mathrm{HCl}$ and $\mathrm{CH}_{3} \mathrm{COONa}$ at infinite dilution are 126.45 , 426.16 and $91.0 \mathrm{Scm}^{2} \mathrm{~mol}^{-1}$ respectively. The molar conductance of $\mathrm{CH}_{3} \mathrm{COOH}$ at infinite dilution is. Choose the right option for your answer.

(1) $201.28 \mathrm{Scm}^{2} \mathrm{~mol}^{-1}$
(2) $390.71 \mathrm{Scm}^{2} \mathrm{~mol}^{-1}$
(3) $698.28 \mathrm{Scm}^{2} \mathrm{~mol}^{-1}$
(4) $540.48 \mathrm{Scm}^{2} \mathrm{~mol}^{-1}$
Ans. (2)
Sol. $\quad \wedge_{\mathrm{m}\left(\mathrm{CH}_{3} \mathrm{COOH}\right)}^{\mathrm{O}}=\wedge_{\mathrm{m}\left(\mathrm{CH}_{3} \mathrm{COONa}\right)}^{\mathrm{O}}+\wedge_{\mathrm{m}(\mathrm{HCl})}^{\mathrm{O}}-\wedge_{\mathrm{m}(\mathrm{NaCl})}^{\mathrm{O}}$

$$
\begin{aligned}
& =91+426.16-126.45 
& =390.71 \mathrm{Scm}^{2} \mathrm{~mol}^{-1}
\end{aligned}
$$

Q78. The $\mathrm{pK}_{\mathrm{b}}$ of dimethylamine and $\mathrm{pK}_{\mathrm{a}}$ of acetic acid are 3.27 and 4.77 respectively at $\mathrm{T}(\mathrm{K})$. The correct option for the pH of dimethylammonium acetate solution is:

(1) 8.50
(2) 5.50
(3) 7.75
(4) 6.25
Ans. (3)
Sol. $\mathbf{p}^{H}=7+\frac{1}{2}\left[\mathbf{p}^{K_{a}}-\mathbf{p}^{K_{b}}\right]$ (Salt of weak acid and weak base)
$=7+\frac{1}{2}[4.77-3.27]$
=7.75

Q79. Match List-I with List-II

\begin{center}
\begin{tabular}{ll}
List-I & List-II 
a. $\mathrm{PCl}_{5}$ & i. Square pyramidal 
b. $\mathrm{SF}_{6}$ & ii. Trigonal planar 
c. $\mathrm{BrF}_{5}$ & iii. Octahedral 
d. $\mathrm{BF}_{3}$ & iv. Trigonal bipyramidal 
\end{tabular}
\end{center}

Choose the correct answer from the options given below.

(1) a - iv; b - iii; c - i; d - ii
(2) a-ii;b-iii;c-iv;d-i
(3) a - iii; b - i; c - iv; d - ii
(4) a-iv;b-iii;c-ii;d-i
Ans. (1)
Sol. Conceptual

Q80. Identify the compound that will react with Hinsberg's reagent to give a solid which dissolves in alkali.
(1) $\\$ IMAGE:https://topnow.s3.us-east-1.amazonaws.com/neet/smile-3edcb589a98c5b736c8a1b550c0c27966dbfa32e.png $\\$
(2) $\\$ IMAGE:https://topnow.s3.us-east-1.amazonaws.com/neet/smile-0abc03c25b23144cdac7c3d3eecc1c897361858c.png $\\$
(3) $\\$ IMAGE:https://topnow.s3.us-east-1.amazonaws.com/neet/smile-a338471d74acc010165bd45083faffc00650c03a.png $\\$
(4) $\\$ IMAGE:https://topnow.s3.us-east-1.amazonaws.com/neet/smile-46295fd04ad2bc0c4922aeff7829414330677721.png $\\$
Ans. (3)
Sol. Only Primary amines reacts with Hinsberg's reagent to give a solid which dissolves in alkali.
$\therefore \mathrm{CH}_{3} \mathrm{CH}_{2} \mathrm{NH}_{2}$

Q81. The right option for the statement "Tyndall effect is exhibited by" is:

(1) NaCl solution
(2) Glucose solution
(3) Starch solution
(4) Urea solution
Ans. (3)
Sol. Conceptual

Q82. $\mathrm{BF}_{3}$ is planar and electron deficient compound. Hybridization and number of electrons around the central atom, respectively are:

(1) $\mathrm{sp}^{3}$ and 4
(2) $\mathrm{sp}^{3}$ and 6
(3) $\mathrm{sp}^{2}$ and 6
(4) $\mathrm{sp}^{2}$ and 8
Ans. (3)
Sol. Conceptual

Q83. A particular station of All India Radio, New Delhi, broad casts on a frequency of $1,368 \mathrm{kHz}$ (kilohertz). The wavelength of the electromagnetic radiation emitted by the transmitter is:
(speed of light, $\mathrm{c}=3.0 \times 10^{8} \mathrm{~ms}^{-1}$ )

(1) 219.3 m
(2) 219.2 m
(3) 2192 m
(4) 21.92 cm
Ans. (1)
Sol. $\quad \mathbf{C}=9 \lambda$
$\lambda=\frac{\mathrm{C}}{9}$
$=\frac{3 \times 10^{8}}{1368 \times 10^{3}}=219.3 \mathrm{~m}$

Q84. Which one among the following is the correct option for right relation between $\mathrm{C}_{\mathrm{p}}$ and $\mathrm{C}_{\mathrm{v}}$ for one mole of ideal gas?

(1) $C_{p}+C_{v}=R$
(2) $C_{p}-C_{v}=R$
(3) $C_{p}=R C_{v}$
(4) $C_{v}=R C_{p}$
Ans. (2)
Sol. $\mathrm{C}_{\mathrm{p}}-\mathrm{C}_{\mathrm{v}}=\mathrm{R}$

Q85. The following solutions were prepared by dissolving 10 g of glucose $\left(\mathrm{C}_{6} \mathrm{H}_{12} \mathrm{O}_{6}\right)$ in 250 ml of water $\left(\mathrm{P}_{1}\right), 10 \mathrm{~g}$ of urea $\left(\mathrm{CH}_{4} \mathrm{~N}_{2} \mathrm{O}\right)$ in 250 ml water ( $\mathrm{P}_{2}$ ) and 10 g of sucrose ( $\mathrm{C}_{12} \mathrm{H}_{22} \mathrm{O}_{11}$ ) in 250 ml of water $\left(\mathrm{P}_{3}\right)$. The right option for the decreasing order of osmotic pressure of these solutions is:

(1) $P_{2}>P_{1}>P_{3}$
(2) $\mathrm{P}_{1}>\mathrm{P}_{2}>\mathrm{P}_{3}$
(3) $P_{2}>P_{3}>P_{1}$
(4) $P_{3}>P_{1}>P_{2}$
Ans. (1) 
Sol. $\pi=$ CRT
$\pi=\frac{\mathrm{W}}{\mathrm{GMW}} \times \frac{1000}{\mathrm{~V}_{(\mathrm{ml})}} \times \mathrm{R} \times \mathrm{T}$
$\pi \propto \frac{1}{\text { GMW }}$
$\mathrm{C}_{6} \mathrm{H}_{12} \mathrm{O}_{6} \rightarrow 180$
$\mathrm{CH}_{4} \mathrm{~N}_{2} \mathrm{O} \rightarrow 60$
$\mathrm{C}_{12} \mathrm{H}_{22} \mathrm{O}_{11} \rightarrow 342$
$\therefore \pi_{\text {urea }}>\pi_{\text {glu cose }}>\pi_{\text {sucrose }}$
$\mathrm{P}_{2}>\mathrm{P}_{1}>\mathrm{P}_{3}$


Q86. The intermediate compound ' X ' in the following chemical reaction is:
$\\$ IMAGE:https://topnow.s3.us-east-1.amazonaws.com/neet/2025_10_11_d2c601cc51a969098abag-14(1).jpg $\\$

1. $\\$ IMAGE:https://topnow.s3.us-east-1.amazonaws.com/neet/smile-9481b9629dfd2e1aa18969a036b88635658bf063.png $\\$

2. $\\$ IMAGE:https://topnow.s3.us-east-1.amazonaws.com/neet/smile-919380ebe0f21ea01f32998763cafc77c304630a.png $\\$

3. $\\$ IMAGE:https://topnow.s3.us-east-1.amazonaws.com/neet/smile-de16b8364e2a2ede0e73fc29a0794de0464625e5.png $\\$

4. $\\$ IMAGE:https://topnow.s3.us-east-1.amazonaws.com/neet/smile-9c4bad59e0c132e923b67115e7742b2404c92eda.png $\\$
Ans. (1)
Sol .$\\$ IMAGE:https://topnow.s3.us-east-1.amazonaws.com/neet/2025_10_11_d2c601cc51a969098abag-14(2).jpg $\\$

Q87. For irreversible expansion of an ideal gas under isothermal condition, the correct option is:

(1) $\Delta \mathrm{U}=0, \Delta \mathrm{~S}_{\text {total }}=0$
(2) $\Delta \mathrm{U} \neq 0, \Delta \mathrm{~S}_{\text {total }} \neq 0$
(3) $\Delta \mathrm{U}=0, \Delta \mathrm{~S}_{\text {total }} \neq 0$
(4) $\Delta \mathrm{U} \neq 0, \Delta \mathrm{~S}_{\text {total }}=0$
Ans. (3)
Sol. For isothermal process $\Delta \mathrm{U}=0$
For irreversible process $\Delta \mathrm{S}_{\text {total }} \neq 0$

Q88. Choose the correct option for the total pressure (in atm ) in a mixture of $4 \mathrm{~g} \mathrm{O}_{2}$ and $2 \mathrm{~g} \mathrm{H}_{2}$ confined in a total volume of one litre at $0^{\circ} \mathrm{C}$ is: (Given : $\mathrm{R}=0.082 \mathrm{~L} \mathrm{~atm} \mathrm{~mol}^{-1} \mathrm{~K}^{-1}, \mathrm{~T}=273 \mathrm{~K}$ )

(1) 2.518
(2) 2.602
(3) 25.18
(4) 26.02
Ans. (3)
Sol. $\mathrm{PV}=\left(\mathrm{n}_{1}+\mathrm{n}_{2}\right) \mathrm{RT}$
$\mathrm{P}=\frac{\left(\mathrm{n}_{1}+\mathrm{n}_{2}\right) \mathrm{RT}}{\mathrm{V}}$
$=\frac{\left(\frac{4}{32}+\frac{2}{2}\right) \times 0.082 \times 273}{1}$
$=25.18$

Q89. The correct option for the value of vapour pressure of a solution at $45^{\circ} \mathrm{C}$ with benzene to octane in molar ratio $3: 2$ is:
[ At $45^{\circ} \mathrm{C}$ vapour pressure of benzene is 280 mm Hg and that of octane is 420 mm Hg . Assume ideal gas)

(1) 160 mm of Hg
(2) 168 mm of Hg
(3) 336 mm of Hg
(4) 350 mm of Hg
Ans. (3)
Sol. $\quad P_{\text {total }}=P_{1}^{0} X_{1}+P_{2}^{0} X_{2}$

$$
\begin{aligned}
& =280 \times \frac{3}{5}+420 \times \frac{2}{5} 
& =336 \mathrm{~mm} \text { of } \mathrm{Hg}
\end{aligned}
$$

Q90. The product formed in the following chemical reaction is:
$\\$ IMAGE:https://topnow.s3.us-east-1.amazonaws.com/neet/2025_10_11_d2c601cc51a969098abag-14(3).jpg $\\$

1. $\\$ IMAGE:https://topnow.s3.us-east-1.amazonaws.com/neet/smile-b05c84bad331ce3d6accab69af257939b340dc4e.png $\\$

2. $\\$ IMAGE:https://topnow.s3.us-east-1.amazonaws.com/neet/smile-e52477b015c9ed393429bce07117ff59e156f00e.png $\\$

3. $\\$ IMAGE:https://topnow.s3.us-east-1.amazonaws.com/neet/smile-634146a0612d88a7bd6859a614cd3ea1456d014a.png $\\$

4. $\\$ IMAGE:https://topnow.s3.us-east-1.amazonaws.com/neet/smile-97d58e6dcc7c153e1d4d303b5ce69f5aa780a7e6.png $\\$
Ans. (4)

Sol. $\\$ IMAGE:https://topnow.s3.us-east-1.amazonaws.com/neet/2025_10_11_d2c601cc51a969098abag-14.jpg $\\$
$\mathrm{NaBH}_{4}$ - reduces only keto group but not ester group.

Q91. Which of the following molecules is non-polar in nature?

(1) $\mathrm{POCl}_{3}$
(2) $\mathrm{CH}_{2} \mathrm{O}$
(3) $\mathrm{SbCl}_{5}$
(4) $\mathrm{NO}_{2}$
Ans. (3)
Sol. Conceptual

Q92. From the following pairs of ions which one is not an iso-electronic pair ?

(1) $\mathrm{O}^{2-}, \mathrm{F}^{-}$
(2) $\mathrm{Na}^{+}, \mathrm{Mg}^{2+}$
(3) $\mathrm{Mn}^{2+}, \mathrm{Fe}^{3+}$
(4) $\mathrm{Fe}^{2+}, \mathrm{Mn}^{2+}$
Ans. (4)
Sol. Conceptual

Q93. The molar conductivity of 0.007 M acetic acid is $20 \mathrm{~S} \mathrm{~cm}^{2} \mathrm{~mol}^{-1}$. What is the dissociation constant of acetic acid ? Choose the correct option.
$\left[\begin{array}{c}\Lambda_{\mathrm{H}^{+}}^{\circ}=350 \mathrm{~S} \mathrm{~cm}^{2} \mathrm{~mol}^{-1} $\\$ \Lambda_{\mathrm{CH}_{3} \mathrm{COO}^{-}}^{\circ}=50 \mathrm{~S} \mathrm{~cm}^{2} \mathrm{~mol}^{-1}\end{array}\right]$

(1) $1.75 \times 10^{-4} \mathrm{~mol} \mathrm{~L}^{-1}$
(2) $2.50 \times 10^{-4} \mathrm{~mol} \mathrm{~L}^{-1}$
(3) $1.75 \times 10^{-5} \mathrm{~mol} \mathrm{~L}^{-1}$
(4) $2.50 \times 10^{-5} \mathrm{~mol} \mathrm{~L}^{-1}$
Ans. (3)
Sol. $\wedge_{\mathrm{m}}=20 ; \wedge_{\mathrm{m}}^{\mathrm{O}}=\lambda_{\mathrm{H}}^{\mathrm{O}}+\lambda_{\mathrm{CH}_{3} \mathrm{COO}^{-}}^{\mathrm{O}}$
$=350+50=400$
$\alpha=\frac{\wedge_{\mathrm{m}}}{\wedge_{\mathrm{m}}^{0}}=\frac{20}{400}=5 \times 10^{-2}$
$\mathrm{K}_{\mathrm{a}}=\mathrm{C} \alpha^{2}=0.007 \times\left(5 \times 10^{-2}\right)^{2}$
$=7 \times 10^{-3} \times 25 \times 10^{-4}$
$=1.75 \times 10^{-5} \mathrm{molL}^{-1}$

Q94. The slope of Arrhenius Plot $\left(\ln \mathrm{kv} / \mathrm{s} \frac{1}{\mathrm{~T}}\right)$ of first order reaction is $-5 \times 10^{3} \mathrm{~K}$. The value of $\mathrm{E}_{\mathrm{a}}$ of the reaction is. Choose the correct option for your answer. 
[Given $\mathrm{R}=8.314 \mathrm{JK}^{-1} \mathrm{~mol}^{-1}$ ]

(1) $41.5 \mathrm{~kJ} \mathrm{~mol}^{-1}$
(2) $83.0 \mathrm{~kJ} \mathrm{~mol}^{-1}$
(3) $166 \mathrm{~kJ} \mathrm{~mol}^{-1}$
(4) $-83 \mathrm{~kJ} \mathrm{~mol}^{-1}$
Ans. (1)
Sol. $\ln \mathrm{k}=\ln \mathrm{A}-\frac{\mathrm{E}_{\mathrm{a}}}{\mathrm{R}} \times \frac{1}{\mathrm{~T}}$
y c m x
slope $=-\frac{E_{a}}{R}$
$-5 \times 10^{3}=-\frac{\mathrm{E}_{a}}{8.314 \times 10^{-3}}$
$\mathrm{E}_{\mathrm{a}}=5 \times 10^{3} \times 8.314 \times 10^{-3}$
$=41.5 \mathrm{kJmol}^{-1}$

Q95. The reagent ' $R$ ' in the given sequence of chemical reaction is :
$\\$ IMAGE:https://topnow.s3.us-east-1.amazonaws.com/neet/2025_10_11_d2c601cc51a969098abag-15(2).jpg $\\$

(1) $\mathrm{H}_{2} \mathrm{O}$
(2) $\mathrm{CH}_{3} \mathrm{CH}_{2} \mathrm{OH}$
(3) HI
(4) $\mathrm{CuCN} / \mathrm{KCN}$
Ans. (2)
Sol. $\mathrm{R}=\mathrm{CH}_{3} \mathrm{CH}_{2} \mathrm{OH}$

Q96. Match List - I with List - II

\begin{center}
\begin{tabular}{|l|l|}
\hline
List-I & List-II 
\hline
\begin{tabular}{l}
(a) CuCl 
\end{tabular} & (i) Hell-VolhardZelinsky reaction 
\hline
\begin{tabular}{l}
(b) 
$\mathrm{NaOX} \rightarrow$ 
\end{tabular} & (ii) Gattermann Koch reaction 
\hline
\begin{tabular}{l}
$\mathrm{R}-\mathrm{CH}_{2}-\mathrm{OH}+$ 
(c) $\mathrm{R}^{\prime} \mathrm{COOH}$ \( \xrightarrow{\text { Conc. } \mathrm{H}_{2} \mathrm{SO}_{4}} \) 
\end{tabular} & (iii) Halofrom reaction 
\hline
(d) \( \begin{aligned} & \mathrm{R}-\mathrm{CH}_{2} \mathrm{COOH} $\\$ & \xrightarrow[\text { (ii) } \mathrm{H}_{2} \mathrm{O}]{\text { (i) } \mathrm{X}_{2} / \operatorname{RedP}} \end{aligned} \) & (iv) Esterificaiton 
\hline
\end{tabular}
\end{center}

(1) a-iv, b-i, c-ii, d-iii
(2) a-iii, b-ii, c-i, d-iv
(3) a-i, b-iv, c-iii, d-ii
(4) a-ii, b-iii, c-iv, d-i
Ans. (4)
Sol. Conceptual

Q97. Match List-I with List - II.

\begin{center}
\begin{tabular}{|l|l|}
\hline
List-I & List-II 
\hline
(a) & (i) Acid rain 
\hline
(b) \( \mathrm{HOCl}(\mathrm{~g}) \xrightarrow{\mathrm{h} v} \) \( \dot{\mathrm{O}} \mathrm{H}+\dot{\mathrm{C}} 1 \) & (ii) Smog 
\hline
(c) & (iii) Ozone depletion 
\hline
(d) & (iv) Tropospheric pollution 
\hline
\end{tabular}
\end{center}

Choose the correct answer from the options given below.

(1) a-i, b-ii, c-iii, d-iv
(2) a-ii, b-iii, c-iv, d-I
(3) a-iv, b-iii, c-i, d-ii
(4) a-iii, b-ii, c-iv, d-i
Ans. (3)
Sol. Conceptual

Q98.
$\mathrm{CH}_{3} \mathrm{CH}_{2} \mathrm{COO}^{-} \mathrm{Na}^{+} \xrightarrow[\text { Heat }]{\mathrm{NaOH},+?} \mathrm{CH}_{3} \mathrm{CH}_{3}+\mathrm{Na}_{2} \mathrm{CO}_{3}$
Consider the above reaction and Identify the missing reagent / chemical.

(1) $\mathrm{B}_{2} \mathrm{H}_{6}$
(2) Red Phosphorus
(3) CaO
(4) DIBALH
Ans. (3)
Sol. $\\$ IMAGE:https://topnow.s3.us-east-1.amazonaws.com/neet/2025_10_11_d2c601cc51a969098abag-15(1).jpg $\\$
$\\$ IMAGE:https://topnow.s3.us-east-1.amazonaws.com/neet/2025_10_11_d2c601cc51a969098abag-15.jpg $\\$
$\mathrm{CH}_{3}-\mathrm{CH}_{3}+\mathrm{Na}_{2} \mathrm{CO}_{3}$

Q99. Match List - I with List-II.

List-I
(a) $\left[\mathrm{Fe}(\mathrm{CN})_{6}\right]^{3-}$
(b) $\left[\mathrm{Fe}\left(\mathrm{H}_{2} \mathrm{O}\right)_{6}\right]^{3+}$
(ii) 0 BM
(iii) 4.90 BM
(iv) 1.73 BM
(c) $\left[\mathrm{Fe}(\mathrm{CN})_{6}\right]^{4-}$
(d) $\left[\mathrm{Fe}\left(\mathrm{H}_{2} \mathrm{O}\right)_{6}\right]^{2+}$
(i) 5.92 BM

List-II

Choose the correct answer from the options given below.

(1) a-iv, b-ii, c-i, d-iii
(2) a-ii, b-iv, c-iii, d-i
(3) a-i, b-iii, c-iv, d-ii
(4) a-iv, b-i, c-ii, d-iii
Ans. (4)
Sol. Conceptual

Q100. In which one of the following arrangements the given sequence is not strictly according to the properties indicated against it ?

(1) $\mathrm{HF}<\mathrm{HCl}<\mathrm{HBr}<\mathrm{HI}$ : Increasing acidic strength
(2) $\mathrm{H}_{2} \mathrm{O}<\mathrm{H}_{2} \mathrm{~S}<\mathrm{H}_{2} \mathrm{Se}<\mathrm{H}_{2} \mathrm{Te}$ : Increasing $\mathrm{pK}_{\mathrm{a}}$ values
(3) $\mathrm{NH}_{3}<\mathrm{PH}_{3}<\mathrm{AsH}_{3}<\mathrm{SbH}_{3}$ : Increasing acidic character
(4) $\mathrm{CO}_{2}<\mathrm{SiO}_{2}<\mathrm{SnO}_{2}<\mathrm{PbO}_{2}$ : Increasing oxidizing power
Ans. (2)
Sol. Conceptual

\section*{BIOLOGY}

Q101.Match List-I with List-II

\begin{center}
\begin{tabular}{|l|l|l|l|}
\hline
\multicolumn{2}{|c|}{List-I} & \multicolumn{2}{c|}{List-II} 
\hline
(a) & Protoplast fusion & (i) & Totipotency 
\hline
(b) & \begin{tabular}{l}
Plant tissue 
culture 
\end{tabular} & (ii) & Pomato 
\hline
(c) & Meristem culture & (iii) & Somaclones 
\hline
(d) & Micropropagation & (iv) & \begin{tabular}{l}
Virus free 
plants 
\end{tabular} 
\hline
\end{tabular}
\end{center}

Choose the correct answer from the options given below

\begin{center}
\begin{tabular}{llllr}
 & (a) & (b) & (c) & (d) 
$(1)$ & (iii) & (iv) & (ii) & (i) 
$(2)$ & (ii) & (i) & (iv) & (iii) 
$(3)$ & (iii) & (iv) & (i) & (ii) 
$(4)$ & (iv) & (iii) & (ii) & (i) 
\end{tabular}
\end{center}
Ans. (2)

Q102. In the equation GPP-R=NPP.

Represents

(1) Radiant energy
(2) Retardation factor
(3) Environment factor
(4) Respiration losses
Ans. (4)

Q103. Which of the following are not secondary metabolites in plants?

(1) Morphine, codeine
(2) Amino acids, Glucose
(3) Vinblastin, Curcumin
(4) Rubber, Gums
Ans. (2)

Q104. The factor that leads to Founder effect in a population is:

(1) Natural selection
(2) Genetic recombination
(3) Mutation
(4) Genetic drift
Ans. (4)

Q105. Amensalism can be represented as:

(1) Species $A(-)$; Species $B(0)$
(2) Species $A(+)$; Species $B(+)$
(3) Species $A(-)$; Species $B(-)$
(4) Species $A(+)$; Species $B(0)$
Ans. (1)

Q106. A typical angiosperm embryo sac at maturity is :

(1) 8-nucleate and 7-celled
(2) 7-nucleate and 8-celled
(3) 7-nucleate and 7-celled
(4) 8-nucleate and 8-celled
Ans. (1)

Q107. During the purification process for recombinant DNA technology, addition of chilled ethanol precipitates out:

(1) RNA
(2) DNA
(3) Histones
(4) Polysaccharids
Ans. (2)

Q108. Gemmae are present in:

(1) Mosses
(2) Pteridophytes
(3) Some Gymnosperms
(4) Some Liverworts
Ans.(4)

Q109. Which of the following stages of meiosis involves division of centromere?

(1) Metaphase I
(2) Metaphase II
(3) Anaphase II
(4) Telophase II
Ans. (3)

Q110. Match List-I with List-II

\begin{center}
\begin{tabular}{|l|l|l|l|}
\hline
\multicolumn{2}{|c|}{List-I} & \multicolumn{2}{c|}{List-II} 
\hline
(a) & Lenticels & (i) & Phellogen 
\hline
(b) & Cork cambium & (ii) & \begin{tabular}{l}
Suberin 
deposition 
\end{tabular} 
\hline
(c) & Secondary cortex & (iii) & \begin{tabular}{l}
Exchange of 
gases 
\end{tabular} 
\hline
(d) & Cork & (iv) & Phelloderm 
\hline
\end{tabular}
\end{center}

Choose the correct answer from the options given below

\begin{center}
\begin{tabular}{lllll}
 & (a) & (b) & (c) & (d) 
$(1)$ & (iv) & (i) & (iii) & (ii) 
$(2)$ & (iii) & (i) & (iv) & (ii) 
$(3)$ & (ii) & (iii) & (iv) & (i) 
$(4)$ & (iv) & (ii) & (i) & (iii) 
\end{tabular}
\end{center}
Ans. (2)

Q111. Plants follow different pathways in response to environment or phases of life to form different kinds of structures. This ability is called:

(1) Elasticity
(2) Flexibility
(3) Plasticity
(4) Maturity
Ans. (3)

Q112. The term used of transfer of pollen grains from anthers of one plant to stigma of a different pant which, during pollination, brings genetically different types of pollen grains to stigma, is:

(1) Xenogamy
(2) Geitonogamy
(3) Chasmogamy
(4) Cleistogamy
Ans. (1)

Q113. Which of the following plants is monoecious?

(1) Carica papaya
(2) Chara
(3) Marchantia polymorpha
(4) Cyas circinalis
Ans. (2)

Q114. Complete the flow chart on central dogma
(a) CDNA $\xrightarrow{\text { (b) }}$ mRNA $\xrightarrow{\text { (c) }}$ (d)

(1) (a)-Replication; (b)-Transcription; (c)-Transduction; (d)-Protein

(2) (a)-Translation; (b)-Replication
(c)-Transcription; (d)-Transduction

(3) (a)-Replication; (b)-Transcription
(c)-Translation; (d)-Protein

(4) (a)-Transduction; (b)-Translation
(c)-Replication; (d)-Protein
Ans. (3)

Q115. Match List-I with List-II

\begin{center}
\begin{tabular}{|l|l|l|l|}
\hline
\multicolumn{2}{|c|}{List-I} & \multicolumn{2}{|c|}{List-II} 
\hline
(a) & Cristae & (i) & Primary constriction in chromosome 
\hline
(b) & Thylnkoids & (ii) & Disc-shaped sacs in Golgi apparatus 
\hline
(c) & Centromere & (iii) & Infoldings in mitochondria 
\hline
(d) & Cisternae & (iv) & Flattened membranous sacs in stroma of plastids 
\hline
\end{tabular}
\end{center}

Choose the correct answer from the options given below

\begin{center}
\begin{tabular}{llllr}
 & (a) & (b) & (c) & (d) 
$(1)$ & (iv) & (iii) & (ii) & (i) 
$(2)$ & (i) & (iv) & (iii) & (ii) 
$(3)$ & (iii) & (iv) & (i) & (ii) 
$(4)$ & (ii) & (iii) & (iv) & (i) 
\end{tabular}
\end{center}
Ans. (3)

Q116. Mutations in plant cells can be induced by:

(1) Kinetin
(2) Infrared rays
(3) Gamma rays
(4) Zeatin
Ans. (3)

Q117. Which of the following statements is not correct?

(1) Pyramid of biomass in sea is generally inverted
(2) Pyramid of biomass in sea is generally upright
(3) Pyramid of energy is always upright
(4) Pyramid of numbers in grassland ecosystem is upright
Ans. (2)


Q118. Inspite of interspecific competition in nature, which mechanism the competing species might have evolved for their survival?

(1) Resource partitioning
(2) Competitive release
(3) Mutualism
(4) Predation
Ans. (1)

Q119. Match List-I with List-II

\begin{center}
\begin{tabular}{|l|l|l|l|}
\hline
\multicolumn{2}{|c|}{List-I} & \multicolumn{2}{c|}{List-II} 
\hline
(a) & Cohesioin & (i) & \begin{tabular}{l}
More attraction in 
liquid phase 
\end{tabular} 
\hline
(b) & Adhesion & (ii) & \begin{tabular}{l}
Mutual attraction 
among water 
molecules 
\end{tabular} 
\hline
(c) & \begin{tabular}{l}
Surface 
tension 
\end{tabular} & (iii) & \begin{tabular}{l}
Water loss in liquid 
phase 
\end{tabular} 
\hline
(d) & Guttation & (iv) & \begin{tabular}{l}
Attraction towards 
polar surfaces 
\end{tabular} 
\hline
\end{tabular}
\end{center}

Choose the correct answer from the options given below

\begin{center}
\begin{tabular}{clccl}
 & (a) & (b) & (c) & (d) 
$(1)$ & (ii) & (iv) & (i) & (iii) 
$(2)$ & (iv) & (iii) & (ii) & (i) 
$(3)$ & (iii) & (i) & (iv) & (ii) 
$(4)$ & (ii) & (i) & (iv) & (iii) 
\end{tabular}
\end{center}
Ans. (1)

Q120. DNA standards on a gel stained with ethidium bromide when viewed under UV radiation, appear as:

(1) Yellow bands
(2) Bright orange bands
(3) Dark red bands
(4) Bright blue bands
Ans. (2)

Q121. Which of the following is an incorrect statement?

(1) Mature sieve tube elements possess a conspicuous nucleus and usual cytoplasmic organelles
(2) Microbodies are present both in plant and animal cells
(3) The perinuclear space forms a barrier between the materials present inside the nucleus and that of the cytoplasm
(4) Nuclear pores act as passages for proteins and RNA molecules in both directions between nucleus and cytoplasm
Ans. (1)

Q122. When gene targeting involving gene amplification is attempted in an individual's tissue to treat disease, it is known as:

(1) Biopiracy
(2) Gene therapy
(3) Molecular diagnosis
(4) Safety testing
Ans. (2)

Q123. Match List-I with List-II

\begin{center}
\begin{tabular}{|l|l|l|l|}
\hline
\multicolumn{2}{|c|}{List-I} & \multicolumn{2}{|c|}{List-II} 
\hline
(a) & Cells with active cell division capacity & (i) & Vascular tissues 
\hline
(b) & Tissues having all cells similar in structure and function & (ii) & Meristematic tissue 
\hline
(c) & Tissue having different types of cells & (iii) & Sclereids 
\hline
(d) & Dead cells with highly thickened walls and narrow Lumen & (iv) & Simple tissue 
\hline
\end{tabular}
\end{center}

Choose the correct answer from the options given below

\begin{center}
\begin{tabular}{llccc}
 & (a) & (b) & (c) & (d) 
(1) & (ii) & (iv) & (i) & (iii) 
(2) & (iv) & (iii) & (ii) & (i) 
(3) & (i) & (ii) & (iii) & (iv) 
(4) & (iii) & (ii) & (iv) & (i) 
\end{tabular}
\end{center}
Ans. (1)

Q124. Which of the following is a correct sequence of steps in a PCR (Polymerase Chain Reaction)?

(1) Denaturation, Annealing, Extension
(2)Denaturation, Extension, Annealing
(3) Extension, Denaturation, Annealing
(4) Annealing, Denaturation, Extension
Ans. (1)

Q125. Which of the following algae produce carrageen

(1) Green algae
(2) Brown algae
(3) Red algae
(4) Blue-green algae
Ans. (3)

Q126. Which of the following is not an application of PCR (Polymerase Chain Reaction)

(1) Molecular diagnosis
(2) Gene amplification
(3) Purification of isolated protein
(4) Detection of gene mutation
Ans. (3)

Q127. Genera like Selaginella and Salvinia produce two kinds of spores. Such plants are known as:

(1) Homosorus
(2) Heterosours
(3) Homosporous
(4) Heterosporous
Ans. (4)

Q128. Diadelphous stamens are found in:

(1) China rose
(3) Citrus
(3) Pea
(4) China rose and citrus
Ans. (3)

Q129. When the centromere is situated in the middle of two equal arms of chromosomes, the chromosome is referred as:

(1) Metacentric
(2) Telocentric
(3) Sub-metacentric
(4) Acrocentric
Ans. (1)

Q130. Which of the following algae contains mannitol as reserve food material?

(1) Ectocarpus
(2) Gracilaria
(3) Voluox
(4) Ulothrix
Ans. (1)

Q131. The amount of nutrients, such as carbon, nitrogen, phosphorus and calcium present in the soil at any given time, is referred as:

(1) Climax
(2) Climax community
(3) Standing state
(4) Standing crop
Ans. (3)

Q132. The first stable product of $\mathrm{CO}_{2}$ fixation in sorghum is:

(1) Pyruvic acid
(2) Oxaloacetic acid
(3) succinic acid
(4) Phosphoglyceric acid
Ans. (2)

Q133. The site of perception of light in plants during photoperiodism is:

(1) Shoot apex
(2) Stem
(3) Axillary bud
(4) Leaf
Ans. (4)

Q134. The plant hormone used to destroy weeds in field is

(1) IAA
(2) NAA
(3) $2,4-\mathrm{D}$
(4) IBA
Ans. (3)

Q135. The production of gametes by the parents, formation of zygotes, the $\mathrm{F}_{1}$ and $\mathrm{F}_{2}$ plants can be understood from a diagram called:

(1) Bullet square
(2) Punch square
(3) Punnett square
(4) Net square
Ans. (3)


Q136. In the exponential growth equation $\mathrm{N}_{\mathrm{t}}=\mathrm{N}_{\mathrm{O}} \mathrm{e}^{\mathrm{rt}} . \mathrm{e}$ represents:

(1) The base of number logarithms
(2) The base of exponential logarithms
(3) The base of natural logarithms
(4) The base of geometric logarithms
Ans. (3)

Q137. Match List-I with List-II

\begin{center}
\begin{tabular}{|l|l|l|lr|}
\hline
\multicolumn{2}{|c|}{List-I} & \multicolumn{3}{c|}{List-II} 
\hline
(a) & Nitrococcus & (i) & Denitrification &  
\hline
(b) & Rhizobium & (ii) & \begin{tabular}{l}
Conversion of
ammonia to nitrite 
\end{tabular} &  
\hline
(c) & Thiobacillus & (iii) & \begin{tabular}{l}
Conversion of 
nitrite to nitrate 
\end{tabular} &  
\hline
(d) & Nitrobacter & (iv) & \begin{tabular}{l}
Conversion of 
atmospheric 
nitrogen 
ammonia 
\end{tabular} & to 
\hline
\end{tabular}
\end{center}

Choose the correct answer from the options given below

\begin{center}
\begin{tabular}{llcrr}
 & (a) & (b) & (c) & (d) 
$(1)$ & (ii) & (iv) & (i) & (iii) 
$(2)$ & (i) & (ii) & (iii) & (iv) 
$(3)$ & (iii) & (i) & (iv) & (ii) 
$(4)$ & (iv) & (iii) & (ii) & (i) 
\end{tabular}
\end{center}
Ans. (1)

Q138. Match List-I with List-II

\begin{center}
\begin{tabular}{|l|l|l|l|}
\hline
\multicolumn{2}{|c|}{List-I} & \multicolumn{2}{|c|}{List-II} 
\hline
(a) & S phase & (i) & Proteins are synthesized 
\hline
(b) & $\mathrm{G}_{2}$ phase & (ii) & Inactive phase 
\hline
(c) & Quiescent stage & (iii) & Interval between mitosis and initiation of DNA replication 
\hline
(d) & $\mathrm{G}_{1}$ phase & (iv) & DNA replication 
\hline
\end{tabular}
\end{center}

Choose the correct answer from the options given below

\begin{center}
\begin{tabular}{llccr}
 & (a) & (b) & (c) & (d) 
$(1)$ & (iii) & (ii) & (i) & (iv) 
$(2)$ & (iv) & (ii) & (iii) & (i) 
$(3)$ & (iv) & (i) & (ii) & (iii) 
$(4)$ & (ii) & (iv) & (iii) & (i) 
\end{tabular}
\end{center}
Ans. (3)

Q139. Identify the correct statement

(1) In capping, methyl guanosine triphosphate is added to the 3 ' end of hnRNA
(2) RNA polymerase binds with Rho factor to terminate the process of transcription in bacteria
(3) The coding strand in a transcription unit is copied to an mRNA
(4) Split gene arrangement is characteristic of prokaryotes
Ans. (2)

Q140. Plasmid pBR322 has PstI restriction enzyme site within gene $\mathrm{amp}^{\mathrm{R}}$ that confers ampicillin resistance. If this enzyme is used for inserting a gene for $\beta$-galactoside production and the recombinant plasmid is inserted in an E.coli strain

(1) It will not be able to confer amplicillin resistance to the host cell
(2) The transformed cells will have the ability to resist ampicillin as well as produce $\beta$ galactoside
(3) It will lead to lysis of host cell
(4) It will be able to produce a novel protein with dual ability
Ans. (1)

Q141. DNA fingerprinting involves identifying differences in some specific regions in DNA sequence, called as:

(1) Satellite DNA
(2) Repetitive DNA
(3) Single nucleotides
(4) Polymorphic DNA
Ans.(2)

Q142. Which of the following statements is correct?

(1) Fusion of two cells is called Karyogamy
(2) Fusion of protoplasms between to motile on non-motile gamets is called plasmogamy
(3) Organisms that depends on livng plants are called saprophytes
(4) Some of the organisms can fix atmospheric nitrogen in specialized cells called sheath cells
Ans. (2)


Q143. Match Column-I with Column-II

\begin{center}
\begin{tabular}{|l|l|}
\hline
Column-I & Column-II 
\hline
(a) $\% \mathrm{~K}_{(5)} \mathrm{C}_{1+2+(2)} \mathrm{A}_{(9)+1} \mathrm{G}_{1}$ & (i) Brassicaceae 
\hline
(b) $\oplus \underset{+}{\chi} \mathrm{K}_{(5)} \widehat{\mathrm{C}_{(5)}} \mathrm{A}_{5} \mathrm{G}_{2}$ & (ii) Liliaceae 
\hline
(c) $\oplus \widehat{\mathrm{P}}_{(3+3)} \mathrm{A}_{3+3} \underline{\mathrm{G}}_{(3)}$ & (iii) Fabaceae 
\hline
\includegraphics[max width=\textwidth]{2025_10_11_d2c601cc51a969098abag-20}
 & (iv) Solanaceae 
\hline
\end{tabular}
\end{center}

Select the correct answer from the options given below.

\begin{center}
\begin{tabular}{lllll}
 & (a) & (b) & (c) & (d) 
(1) & (iii) & (iv) & (ii) & (i) 
(2) & (i) & (ii) & (iii) & (iv) 
(3) & (ii) & (iii) & (iv) & (i) 
(4) & (iv) & (ii) & (i) & (iii) 
\end{tabular}
\end{center}
Ans. (1)

Q144. Now a days it is possible to detect the mutated gene causing cancer by allowing radioactive probe to hybridise its complimentary DNA in a clone of cells. Followed by its detection using autoradiography because.

(1) Mutated gene partially appears on a photographic film
(2) Mutated gene completely and clearly appears on a photographic film
(3) Mutated gene does not appear on a photographic film as the probe has no complimentarity with it
(4) Mutated gene does not appear on photographic film as the probe has complimentarity with it
Ans. (3)

Q145. Which of the flowing statements is incorrect?

() Both ATP and NADPH $+\mathrm{H}^{+}$are synthesized during non-cyclic photophosphorylation
(2) Stroma lamellae have PS I only and lack NADP reductase
(3) Grana lamellae have both PS I only PS II
(4) Cyclic photophosphorylation involves both PS I and PS II
Ans. (4)


Q146. Which of the following statements is incorrect?

(1) During aerobic respiration, role of oxygen is limited to the terminal stage
(2) In ETC (Electron Transport Chain), one molecule of $\mathrm{NADH}+\mathrm{H}^{+}$gives rise to 2 ATP molecules, and one $\mathrm{FADH}_{2}$ gives rise to 3 ATP molecules
(3) ATP is synthesized through complex $V$
(4) Oxidation-reduction reactions produce proton gradient in respiration
Ans. (2)

Q147. Match List-I with List-II

\begin{center}
\begin{tabular}{|l|l|l|l|}
\hline
\multicolumn{2}{|c|}{List-I} & \multicolumn{2}{c|}{List-II} 
\hline
(a) & Protein & (i) & \begin{tabular}{l}
C=C double 
bonds 
\end{tabular} 
\hline
(b) & \begin{tabular}{l}
Unsaturated 
fatty acid 
\end{tabular} & (ii) & \begin{tabular}{l}
Phosphodiester 
bonds 
\end{tabular} 
\hline
(c) & Nucleic acid & (iii) & \begin{tabular}{l}
Glycosidic 
bonds 
\end{tabular} 
\hline
(d) & Polysaccharide & (iv) & Peptide bonds 
\hline
\end{tabular}
\end{center}

Choose the correct answer from the options given below

\begin{center}
\begin{tabular}{llcrr}
 & (a) & (b) & (c) & (d) 
(1)$ & (iv) & (i) & (ii) & (iii) 
(2)$ & (i) & (iv) & (iii) & (ii) 
(3)$ & (ii) & (i) & (iv) & (iii) 
(4)$ & (iv) & (iii) & (i) & (ii) 
\end{tabular}
\end{center}
Ans. (1)

Q148. What is the role of RNA polymerase III in the process of transcription in eukaryotes?

(1) Transcribes rRNAs (28S, 18 S and 5.8 s )
(2) Transcribes tRNA, 5 s rRNA and snRNA
(3) Transcribes precursor of mRNA
(4) Transcribes only snRNAs
Ans. (2)

Q149. Select the correct pair

\begin{center}
\begin{tabular}{|l|l|}
\hline
(1) Large colorless empty cells in the epidermis of grass leaves & Subsidary cells 
\hline
(2) In dicot leaves, vascular - bundles are surrounded by large thick-walled cells & Conjunctive tissue 
\hline
(3) Cells of medullary rays that form part of cambial ring & Interfascicular cambium 
\hline
(4) Loose parenchyma cells Rupturing the epidermis and forming a lens-shaped opening in bark & Spongy parenchyma 
\hline
\end{tabular}
\end{center}
Ans. (3)

Q150. In some members of which of the following pairs of families, pollen grains retain their viability for months after release?

(1) Poaceae; Rosaceae
(2) Poaceae; Leguminosae
(3) Poaceae; Solanaceae
(4) Rosacea; Leguminosae
Ans. (4)

Q151. Match List-I with List-II

\begin{center}
\begin{tabular}{|l|l|}
\hline
\multicolumn{1}{|c|}{List-I} & \multicolumn{1}{|c|}{List-II} 
\hline
(a) Vaults & \begin{tabular}{l}
(i) Entry of sperm through 
Cervix is blocked 
\end{tabular} 
\hline
(b) IUDs & (ii) Removal of Vas deferens 
\hline
(c) Vasectomy & \begin{tabular}{l}
(iii) Phagocytosis of sperms 
within the Uterus 
\end{tabular} 
\hline
\begin{tabular}{l}
(d) 
Tubectomy 
\end{tabular} & \begin{tabular}{l}
(iv) Removal of fallopian 
tube 
\end{tabular} 
\hline
\end{tabular}
\end{center}

Choose the correct answer form the option given below

(1) a-iv, b-ii, c-i, d-iii
(2) a-i, b-iii, c-ii, d-iv
(3) a-ii, b-iv, c-iii, d-i
(4) a-iii, b-i, c-iv, d-ii
Ans. (2)

Q152. Which of the following statements wrongly represents the nature of smooth muscle?

(1) These muscle have no striations
(2) They are involuntary muscles
(3) Communication among the cells is performed by intercalated discs
(4) These muscles are present in the wall of blood vessels
Ans. (3)

Q153. The organelles that are included in the endomembrane system are

(1) Endoplasmic reticulum, Mitochondria, Ribosomes and Lysosomes
(2) Endoplasmic reticulum, Golgi complex, Lysosomes and Vacuoles
(3) Golgi complex, Mitochondria, Ribosomes and Lysosomes
(4) Golgi complex, Endoplasmic reticulum, Mitochondria and Lysosomes
Ans. (2)

Q154. Succus entericus is referred to as

(1) Pancreatic juice
(2) Intestinal juice
(3) Gastric juice
(4) Chyme
Ans. (2)

Q155. Which one of the following is an example of Hormone releasing IUD?

(1) CuT
(2) LNG 20
(3) Cu 7
(4) Multiload 375
Ans. (2)

Q156. Which one of the following belongs to the family Muscidae?

(1) Fire fly
(2) Grasshopper
(3) Cockroach
(4) House fly
Ans. (4)

Q157. If Adenine makes $30 \%$ of the DNA molecule, what will be the percentage of Thymine, Guanine and Cytosine in it?

(1) T:20; G:30; C:20
(2) T:20; G:20; C:30
(3) T:30; G:20; C:20
(4) $\mathrm{T}: 20 ; \mathrm{G}: 25 ; \mathrm{C}: 25$
Ans. (3)

Q158. Receptors of sperm binding in mammals are present on

(1) Corona radiata
(2) Vitelline membrane
(3) Perivitelline space
(4) Zona pellucida
Ans. (4)

Q159. Which of the following is not an objective of Biofortification in crops?

(1) Improve protein content
(2) Improve resistance to diseases
(3) Improve vitamin content
(4) Improve micronutrient and mineral content
Ans. (2)

Q160. The centriole undergoes duplication during

(1) S-phase
(2) Prophase
(3) Metaphase
(4) $G_{2}$ phase
Ans. (1)

Q161. Chronic auto immune disorder affecting neuro muscular junction leading to fatigue, weakening and paralysis of skeletal muscle is called as

(1) Arthritis
(2) Muscular dystrophy
(3) Myasthenia gravis
(4) Gout
Ans. (3)

Q162. Match List-I with List-II

\begin{center}
\begin{tabular}{|l|l|}
\hline
\multicolumn{1}{|c|}{List-I} & \multicolumn{1}{c|}{List-II} 
\hline
(a) Metamerism & (i) Coelenterata 
\hline
(b) Canal system & (ii) Ctenophora 
\hline
(c) Comb plates & (iii) annelida 
\hline
(d) Cnidoblasts & (iv) Porifera 
\hline
\end{tabular}
\end{center}

Choose the correct answer form the option given below

(1) a-iv, b-iii, c-i, d-ii
(2) a-iii, b-iv, c-i, d-ii
(3) a-iii, b-iv, c-ii, d-i
(4) a-iv, b-i, c-ii, d-iii
Ans. (3)

Q163. During the process of gene amplification using PCR, if very high temperature is not maintained in the beginning, then which of the following steps of PCR will be affected first?

(1) Annealing
(2) Extension
(3) Denaturation
(4) Ligation
Ans. (3)

Q164. Read the following statements
(a) Metagenesis is observed in Helminths
(b) Echinoderms are triploblastic and coelomate animals
(c) Round worms have organ-system level of body organization
(d) Comb plates present in ctenophores help in digestion
(e) Water vascular system is characteristic of Echinoderms
Choose the correct answer form the option given below

(1) c, d and e are correct
(2) a, b and c are correct
(3) a, d and e are correct
(4) b, c and e are correct
Ans. (4)

Q165. Dobson units are used measure thickness of

(1) CFCs
(2) Stratosphere
(3) Ozone
(4) Troposphere
Ans. (3)

Q166. Which is the 'Only enzyme'' that has "Capability" to a catalyse Initiation. Elongation and Termination the process of transcription in prokaryotes?

(1) DNA dependent DNA polymerase
(2) DNA dependent RNA polymerase
(3) DNA ligase
(4) DNase
Ans. (2)

Q167. A specific recognition sequence identified by endonucleases to make cuts at specific positions within the DN is

(1) Degenerate primer sequence
(2) Okazaki sequences
(3) Palindromic Nucleotide sequences
(4) Poly (A) tail sequences
Ans. (3)

Q168. The fruit fly has 8 chromosome ( 2 n ) in each cell. During interphase of Mitosis if the number of chromosomes at $\mathrm{G}_{1}$ phase is 8 , what would be the number of chromosomes after S phase?

(1) 8
(2) 16
(3) 4
(4) 32
Ans. (1)

Q169. Sphincter of oddi is present at

(1) Ileo-caecal junction
(2) Junction of hepato-pancreatic duct and duodenum
(3) Gastro-oesophageal junction
(4) Junction of jejunum and duodenum
Ans. (2)

Q170. Math List-I with List-II

\begin{center}
\begin{tabular}{|l|l|}
\hline
\multicolumn{1}{|c|}{List-I} & \multicolumn{1}{c|}{List-II} 
\hline
(a) Aspergillus niger & (i) Acetic Acid 
\hline
(b) Acetobacter aceti & (ii) Lactic Acid 
\hline
\begin{tabular}{l}
(c) Clostridium 
butylicum 
\end{tabular} & (iii) Citric Acid 
\hline
(d) Lacto bacillus & (iv) Butyric Acid 
\hline
\end{tabular}
\end{center}

Choose the correct answer form the option given below

(1) a-iii, b-i, c-iv, d-ii
(2) a-i, b-ii, c-iii, d-iv
(3) a-ii, b-iii, c-i, d-iv
(4) a-iv, b-ii, c-i, d-iii
Ans. (1)

Q171. Which one of the following organisms bears hallow and pneumatic long bones?

(1) Neophron
(2) Hemidactylus
(3) Macropus
(4) Ornithorhynchus
Ans. (1)

Q172. In a cross between a male and female, both heterozygous for sickle cell anaemia gene, what percentage of the progeny will be diseased?

(1) $50 \%$
(2) $75 \%$
(3) $25 \%$
(4) $100 \%$
Ans. (3)

Q173. The partial pressures (in mm Hg ) of oxygen $\left(\mathrm{O}_{2}\right)$ and carbon dioxide $\left(\mathrm{CO}_{2}\right)$ at alveoli ( the site of diffusion) are

(1) $\mathrm{pO}_{2}=104$ and $\mathrm{pCO}_{2}=40$
(2) $\mathrm{pO}_{2}=40$ and $\mathrm{pCO}_{2}=45$
(3) $\mathrm{pO}_{2}=95$ and $\mathrm{pCO}_{2}=40$
(4) $\mathrm{pO}_{2}=159$ and $\mathrm{pCO}_{2}=0.3$
Ans. (1)

Q174. Veneral diseases can spread through
(a) Using sterile needles
(b) Transfusion of blood from infected person
(c) Infected mother to foetus
(d) Kissing
(e) Inheritance

Choose the correct answer form the option given below

(1) a, b and c only
(2) b, c and d only
(3) b and c only
(4) a and c only
Ans. (3)

Q175. Which of the following RNAs is not required for the synthesis of protein?

(1) mRNA
(2) tRNA
(3) rRNA
(4) siRNA
Ans. (4)

Q176. Select the favourable conditions required for the formation of oxyhaemoglobin at the alveoli

(1) High $\mathrm{pO}_{2}$, low $\mathrm{pCO}_{2}$, less $\mathrm{H}^{+}$, lower temperature
(2) Low $\mathrm{pO}_{2}$, high $\mathrm{pCO}_{2}$, more $\mathrm{H}^{+}$, higher temperature
(3) High $\mathrm{pO}_{2}$, high $\mathrm{pCO}_{2}$, less $\mathrm{H}^{+}$, higher temperature
(4) Low $\mathrm{pO}_{2}$, low $\mathrm{pCO}_{2}$, more $\mathrm{H}^{+}$, high temperature
Ans. 1
Ans. (1)

Q177. Match the following

\begin{center}
\begin{tabular}{|l|l|}
\hline
\multicolumn{1}{|c|}{List-I} & \multicolumn{1}{c|}{List-II} 
\hline
(a) Physalia & (i) Pearly oyster 
\hline
(b) Limulus & \begin{tabular}{l}
(ii) Portuguese Man 
of War 
\end{tabular} 
\hline
(c) Ancylostoma & (iii) Living fossile 
\hline
(d) Pinctada & (iv) Hook worm 
\hline
\end{tabular}
\end{center}

Choose the correct answer form the option given below

(1) a-ii, b-iii, c-i, d-iv
(2) a-iv, b-i, c-iii, d-ii
(3) a-ii, b-iii, c-iv, d-i
(4) a-i, b-iv, c-iii, d-ii
Ans. (3)

Q178. Which enzyme is responsible for the conversion of inactive fibrinogens to fibrins?

(1) Thrombin
(2) Renin
(3) Epinephrine
(4) Thrombokinase
Ans. (1)

Q179. For effective treatment of the disease, early diagnosis and understanding its pathophysiology is very important. Which of the following molecular diagnostic techniques is very useful for early detection?

(1) Western Blotting Technique
(2) Southern Blotting Technique
(3) ELISA Technique
(4) Hybridization Technique
Ans. (3)

Q180. Identify the incorrect pair

(1) Alkaloids - Codeine
(2) Toxin - Abrin
(3) Lectins - Concanavalin A
(4) Drugs - Ricin
Ans. (4)

Q181. Which stage of meiotic prophase shows terminalisation of chiasmata as its distinctive feature?

(1) Leptotene
(2) Zygotene
(3) Diakinesis
(4) Pachytene
Ans. (3)

Q182. With regard to insulin choose correct options
(a) C -peptide is not present in mature insulin
(b) The insulin produced by rDNA technology has C-peptide
(c) The pro-insulin has C -peptide
(d) A-peptide and B-peptide of insulin are interconnected by disulphide bridges
Choose the correct answer from the options given below

(1) b and d only
(2) b and c only
(3) a, c and d only
(4)a and d only
Ans. (3)

Q183. Person with 'AB' blood group are called "Universal recipients". This is due to

(1) Absence of antigens $A$ and $B$ on the surface of RBCs
(2) Absences of antigens $A$ and $B$ in plasma
(3) Presence of antibodies, anti-A and anti-B, on RBCs
(4) Absence of antibodies, anti-A and anti-B in plasma
Ans. (4)

Q184. Erythropoietin hormone which stimulates R.B.C formation is produced by

(1) Alpha cells of pancreas
(2) The cells of rostral adenohypophysis
(3) The cells of bone marrow
(4) Juxtaglomerular cells of the kidney
Ans. (4)

Q185. Which of the following characteristic is incorrect with respect to cockroach?

(1) A ring of gastric caeca is present at the junction of midgut and hind gut
(2) Hypopharynx lies within the cavity enclosed by the mouth parts
(3) In females, $7^{\text {th }}-9^{\text {th }}$ sterna together form a genital pouch
(4) $10^{\text {th }}$ abdominal segment in both sexes, bears a pair of anal cerci
Ans. (1)

Q186. The adenosine deaminase deficiency results into

(1) Dysfunction of immune system
(2) Parkinson's disease
(3) Digestive disorder
(4) Addison's disease
Ans. (1)

Q187. Which of the following is not a step in Multiple Ovulation Embryo Transfer Technology (MOET)?

(1) Cow is administered hormone having LH like activity for super ovulation
(2) Cow yields about 6-8 eggs at a time
(3) Cow is fertilized by artificial insemination
(4) Fertilized eggs are transferred to surrogate mother at 8-32 cell stage
Ans. (1)

Q188. Math List-I with List-II

\begin{center}
\begin{tabular}{|l|l|}
\hline
List-I & List-II 
\hline
(a) Adaptive radiation & (i) Selection of resistant varieties due to excessive use of herbicides and pesticides 
\hline
(b) Convergent evolution & (ii) Bones of forelimbs in Man and Whale 
\hline
(c) Divergent evolution & (iii) Wings of Butterfly and Bird 
\hline
(d) Evolution by anthropogenic action & (iv) Darwin Finches 
\hline
\end{tabular}
\end{center}

Choose the correct answer form the options given below

(1) a-iv, b-iii, c-ii, d-i
(2) a-iii, b-ii, c-i, d-iv
(3) a-ii, b-i, c-iv, d-iii
(4) a-i, b-iv, c-iii, d-ii
Ans. (1)

Q189. Which one of the following statements about Histones is wrong?

(1) Histones are organized to form a unit of 8 molecules
(2) The pH of histones is slightly acidic
(3) Histones are rich in amino acids - Lysine and Arginine
(4) Histones carry positive charge in the side chain
Ans. (2)


Q190. Which of the following secretes the hormone, relaxin, during the later phase of pregnancy?

(1) Graffian follicle
(2) Corpus luteum
(3) Foetus
(4) Uterus
Ans. (2)

Q191. Following are the statements with reference to 'lipids'
(a) Lipids having only single bonds are called unsaturated fatty acids
(b) Lecithin is a phospholipid
(c) Trihydroxy propane is glycerol
(d) Palmitic acid has 20 carbon atoms includes carboxyl carbon
(e) Arachidonic acid has 16 carbon atoms

Choose the correct answer from the options given below

(1) a and b only
(2) c and d only
(3) b and c only
(4) b and e only
Ans. (3)

Q192. Match List-I with List-II

\begin{center}
\begin{tabular}{|l|l|}
\hline
\multicolumn{1}{|c|}{List-I} & \multicolumn{1}{c|}{List-II} 
\hline
(a) Filariasis & \begin{tabular}{l}
(i) Haemophilus 
influenza 
\end{tabular} 
\hline
(b) Amoebiasis & (ii) Trichophyton 
\hline
(c) Pneumonia & \begin{tabular}{l}
(iii) Wuchereria 
bancrofti 
\end{tabular} 
\hline
(d) Ringworm & \begin{tabular}{l}
(iv) Entamoeba 
histolytica 
\end{tabular} 
\hline
\end{tabular}
\end{center}

Choose the correct answer from the options given below

(1) a-iv, b-i, c-iii, d-ii
(2) a-iii, b-iv, c-i, d-ii
(3) a-i, b-ii, c-iv, d-iii
(4) a-ii, b-iii, c-i, d-iv
Ans. (2)

Q193. Identify the types of cell junctions that help to strong the leakage of the substance across a tissues and facilitation of communication with neighbouring cells via rapid transfer of ions and molecules.

(1) Gap junctions and Adhering junctions respectively
(2) Tight junctions and Gap junctions respectively
(3) Adhering junctions and Tight junctions respectively
(4) Adhering junctions and Gap junctions respectively
Ans. (2)

Q194. During muscular contraction which of the following events occur?
(a) 'H' zone disappears
(b) 'A' band widens
(c) 'I' band reduces in width
(d) Myosine hydrolyzes ATP, releasing the ADP and Pi
(e) Z-lines attached to actins are pulled inwards.
Choose the correct answer from the options given below

(1) a, c, d, e only
(2) a, b, c, d only
(3) b, c, d, e only
(4) b, d, e a, only
Ans. (1)

Q195. Following are the statements about prostomium of earthworm
(a) It serves as a covering for mouth
(b) It helps to open cracks in the soil into which it can crawl
(c) It is one of the sensory structures
(d) it is the first body segment

Choose the correct answer from the options given below

(1) a, b, and c are correct
(2) a, b and d are correct
(3) a, b, c and d are correct
(4) b and c are correct
Ans. (1)

Q196. Assertion (A): A person goes to high altitude and experiences 'altitude sickness' with symptoms like breathing difficulty and heart palpitations
Reason (R): Due to low atmospheric pressure at high altitude, the body does not get sufficient oxygen.
In the light of the above statements, choose the correct answer from the options given below

(1) Both $A$ and $R$ are true and $R$ is the correct explanation of $A$
(2) Both $A$ and $R$ are true but $R$ is not the correct explanation of $A$
(3) $A$ is true but $R$ is false
(4) A is false but R is true
Ans. (1)

Q197. Which of these is not an important component of initiation of parturition in humans?

(1) Increase in estrogen and progesterone ratio
(2) Synthesis of prostaglandins
(3) Release of Oxytocin
(4) Release of Prolactin
Ans. (4)

Q198. Match List-I with List-II

\begin{center}
\begin{tabular}{|l|l|}
\hline
\multicolumn{1}{|c|}{List-I} & \multicolumn{1}{c|}{List-II} 
\hline
(a) Allen's Rule & (i) Kangaroo rat 
\hline
\begin{tabular}{l}
(b) Physiological 
adaptation 
\end{tabular} & (ii) Desert lizard 
\hline
\begin{tabular}{l}
(c) Behavioural 
adaptation 
\end{tabular} & \begin{tabular}{l}
(iii) Marine fish at 
depth 
\end{tabular} 
\hline
\begin{tabular}{l}
(d) Biochemical 
adaptation 
\end{tabular} & (iv) Polar seal 
\hline
\end{tabular}
\end{center}

Choose the correct answer from the options given below

(1) a-iv, b-ii, c-iii, d-i
(2) a-iv, b-i, c-iii, d-ii
(3) a-iv, b-i, c-ii, d-iii
(4) a-iv, b-iii, c-ii, d-i
Ans. (3)

Q199. Match List-I with List-II

\begin{center}
\begin{tabular}{|l|l|}
\hline
\multicolumn{1}{|c|}{List-I} & \multicolumn{1}{c|}{List-II} 
\hline
(a) Scapula & (i) Cartilaginous joints 
\hline
(b) Cranium & (ii) Flat bone 
\hline
(c) Sternum & (iii) Fibrous joints 
\hline
\begin{tabular}{l}
(d) Vertebral 
column 
\end{tabular} & (iv) Triangular flat bone 
\hline
\end{tabular}
\end{center}

Choose the correct answer from the options given below

(1) a-i, b-iii, c-ii, d-iv
(2) a-ii, b-iii, c-iv, d-i
(3) a-iv, b-ii, c-iii, d-i
(4) a-iv, b-iii, c-ii, d-i
Ans. (4)

Q200. Statement-I: The codon 'AUG' codes for methionine and phenylalanine
Statement-II: 'AAA' and 'AAG' both codons code for the amino acid lysine
In the light of the above statements, Choose the correct answer from the options given below

(1) Both Statement I and Statement II are true
(2) Both Statement I and Statement II are false
(3) Statement I is correct but Statement II is false
(4) Statement I is incorrect but Statement II is true
Ans. (4)

\end{document}